\documentclass[12pt,a4paper]{article}
\usepackage[T1]{fontenc}
\usepackage[utf8]{inputenc}
\usepackage{lmodern}
\usepackage{amsmath,amssymb,amsthm,mathtools}
\usepackage{geometry}
\geometry{margin=1in}
\usepackage{microtype}
\usepackage{hyperref}
\usepackage{xcolor}
\usepackage{enumitem}
\usepackage{fancyhdr}
\usepackage{bookmark}

\usepackage{centernot}

\hypersetup{
	colorlinks=true,
	linkcolor=blue!50!black,
	urlcolor=blue!50!black,
	citecolor=blue!50!black,
	pdftitle={Theory of Commutative Algebra 7},
	pdfauthor={},
	pdfsubject={Solved problems and theory notes in commutative algebra},
	pdfcreator={OpenAI}
}

\setlist[itemize]{topsep=4pt,itemsep=2pt,parsep=0pt,partopsep=0pt}
\setlength{\parskip}{0.55em}
\setlength{\parindent}{0pt}

\setcounter{secnumdepth}{2}


\setcounter{tocdepth}{2}

\pagestyle{fancy}


\fancyhf{}
\lhead{Theory of Algebraic Topology 1}
\rhead{\thepage}
\renewcommand{\headrulewidth}{0.4pt}

\newtheoremstyle{notespace}{6pt}{6pt}{\normalfont}{} {\bfseries}{.}{0.5em}{}
\theoremstyle{notespace}
\newtheorem{definition}{Definition}[section]
\newtheorem{remark}{Remark}[section]

\DeclareMathOperator{\Spec}{Spec}
\DeclareMathOperator{\Ass}{Ass}
\DeclareMathOperator{\Ann}{Ann}
\DeclareMathOperator{\Supp}{Supp}
\DeclareMathOperator{\Hom}{Hom}
\DeclareMathOperator{\Min}{Min}
\DeclareMathOperator{\Ker}{Ker}
\DeclareMathOperator{\Frac}{Frac}

\DeclareMathOperator{\depth}{depth}
\DeclareMathOperator{\Tor}{Tor}
\DeclareMathOperator{\Ext}{Ext}

\DeclareMathOperator{\coker}{coker}
\DeclareMathOperator{\im}{im}


\title{\vspace{-1.5em}\bfseries Theory of Algebraic Topology 1}
\author{}
\date{\today}

\begin{document}
	\maketitle
	\thispagestyle{plain}
		\begin{center}
		\begin{minipage}{0.92\textwidth}
			Lecture notes with worked examples in algebraic topology
		\end{minipage}
	\end{center}
		\pdfbookmark[1]{Contents}{contents}\tableofcontents
	\clearpage
	
\section{Problem 1 --- The antipodal map on odd-dimensional spheres}

\subsection{Exact statement}

Let
\[
a:S^n\to S^n,\qquad a(x)=-x,
\]
be the antipodal map. Show that \(a\) is homotopic to the identity map when \(n\) is odd.

\subsection{Solution}

We must show that if \(n\) is odd, then there exists a continuous map
\[
H:S^n\times [0,1]\to S^n
\]
such that
\[
H(x,0)=a(x)=-x
\qquad\text{and}\qquad
H(x,1)=x.
\]

\subsubsection{Step 1: The case \(n=1\)}

First consider
\[
S^1=\{z\in \mathbb C:|z|=1\}.
\]
The antipodal map is
\[
a(z)=-z.
\]
On the circle, multiplication by \(-1\) is just rotation by angle \(\pi\). So we can connect \(-1\) to \(1\) by a continuous path in the unit circle.

Define
\[
H(z,t)=e^{i\pi(1-t)}z.
\]
Then \(H\) is continuous, and since \(|e^{i\pi(1-t)}|=1\), we have \(H(z,t)\in S^1\) for all \(z\in S^1\), \(t\in[0,1]\).

Moreover,
\[
H(z,0)=e^{i\pi}z=-z=a(z),
\]
and
\[
H(z,1)=e^0 z=z.
\]
Thus \(a\) is homotopic to the identity on \(S^1\).

\subsubsection{Step 2: The general odd-dimensional case}

Now suppose \(n\) is odd. Then we may write
\[
n=2m+1
\]
for some integer \(m\ge 0\). Realize \(S^{2m+1}\) as the unit sphere in \(\mathbb C^{m+1}\):
\[
S^{2m+1}
=
\left\{
(z_0,\dots,z_m)\in \mathbb C^{m+1}
:
\sum_{j=0}^m |z_j|^2=1
\right\}.
\]

Define
\[
H\big((z_0,\dots,z_m),t\big)
=
\big(e^{i\pi(1-t)}z_0,\dots,e^{i\pi(1-t)}z_m\big).
\]

We check that this is well-defined. Since multiplication by a complex number of modulus \(1\) preserves norms,
\[
\sum_{j=0}^m \left|e^{i\pi(1-t)}z_j\right|^2
=
\sum_{j=0}^m |z_j|^2
=
1.
\]
Hence \(H((z_0,\dots,z_m),t)\in S^{2m+1}\) for all \(t\in[0,1]\).

The map \(H\) is continuous, and at the endpoints we obtain:
\[
H((z_0,\dots,z_m),0)
=
(-z_0,\dots,-z_m)
=
-(z_0,\dots,z_m)
=
a(z),
\]
while
\[
H((z_0,\dots,z_m),1)
=
(z_0,\dots,z_m).
\]

Therefore \(H\) is a homotopy from the antipodal map \(a\) to the identity map on \(S^{2m+1}\).

Thus, whenever \(n\) is odd, the antipodal map
\[
a(x)=-x
\]
is homotopic to the identity on \(S^n\).

\subsection{Conclusion}

If \(n\) is odd, then \(S^n=S^{2m+1}\) can be viewed as the unit sphere in \(\mathbb C^{m+1}\), and the homotopy
\[
H(z,t)=e^{i\pi(1-t)}z
\]
provides a continuous deformation from the antipodal map to the identity. Hence
\[
a\simeq \mathrm{id}_{S^n}
\qquad\text{for odd }n.
\]

\subsection{Remark}

For even \(n\), the antipodal map is not homotopic to the identity. Indeed,
\[
\deg(a)=\deg(-\mathrm{id}_{S^n})=(-1)^{n+1}.
\]
So if \(n\) is even, then \(\deg(a)=-1\), whereas
\[
\deg(\mathrm{id}_{S^n})=1.
\]
Since homotopic maps have the same degree, \(a\) cannot be homotopic to the identity when \(n\) is even.	
	
	
\section{Homotopic maps: definition, intuition, and examples}

\subsection{Exact statement}

Let \(X\) and \(Y\) be topological spaces, and let
\[
f,g:X\to Y
\]
be continuous maps. We say that \(f\) and \(g\) are \emph{homotopic} if there exists a continuous map
\[
H:X\times[0,1]\to Y
\]
such that for every \(x\in X\),
\[
H(x,0)=f(x),\qquad H(x,1)=g(x).
\]

The map \(H\) is called a \emph{homotopy} from \(f\) to \(g\), and we write
\[
f\simeq g.
\]

\subsection{Intuition}

A homotopy is a continuous deformation of one map into another. The parameter
\[
t\in[0,1]
\]
plays the role of time:
\[
t=0 \Rightarrow f,\qquad t=1 \Rightarrow g.
\]

For each fixed \(t\in[0,1]\), the formula
\[
H_t(x)=H(x,t)
\]
defines a continuous map \(H_t:X\to Y\). Thus a homotopy is a continuous family of maps joining \(f\) to \(g\).

Equivalently, for each \(x\in X\), the point \(f(x)\in Y\) moves continuously along a path until it reaches \(g(x)\).

\subsection{Null-homotopic maps}

A map \(f:X\to Y\) is called \emph{null-homotopic} if it is homotopic to a constant map. That is, there exists \(y_0\in Y\) such that
\[
f\simeq c_{y_0},
\qquad
c_{y_0}(x)=y_0.
\]

\subsection{Examples}

\subsubsection{Example 1: Maps \(\mathbb R\to\mathbb R\)}

Let
\[
f(x)=x,\qquad g(x)=2x.
\]
Define
\[
H(x,t)=(1-t)x+t(2x)=(1+t)x.
\]
Then
\[
H(x,0)=x=f(x),
\qquad
H(x,1)=2x=g(x).
\]
Hence \(f\simeq g\).

\subsubsection{Example 2: The identity map on \(\mathbb R^n\) is homotopic to a constant map}

Let
\[
\mathrm{id}_{\mathbb R^n}(x)=x,
\qquad
c(x)=0.
\]
Define
\[
H(x,t)=(1-t)x.
\]
Then
\[
H(x,0)=x,
\qquad
H(x,1)=0.
\]
Thus
\[
\mathrm{id}_{\mathbb R^n}\simeq c.
\]
So \(\mathbb R^n\) can be continuously shrunk to a point.

\subsubsection{Example 3: Straight-line homotopy in a convex set}

Suppose \(Y\subseteq \mathbb R^m\) is convex, and let
\[
f,g:X\to Y
\]
be continuous maps. Define
\[
H(x,t)=(1-t)f(x)+t\,g(x).
\]
Since \(Y\) is convex, \(H(x,t)\in Y\) for all \(x\in X\) and \(t\in[0,1]\). Therefore \(H\) is a homotopy from \(f\) to \(g\).

In particular, any two maps into \(\mathbb R^m\) are homotopic, and every map into \(\mathbb R^m\) is null-homotopic.

\subsubsection{Example 4: Two constant maps}

Let \(X\) be any topological space, and let \(Y\) be path-connected. Let
\[
c_{y_0}(x)=y_0,
\qquad
c_{y_1}(x)=y_1
\]
be two constant maps. Since \(Y\) is path-connected, there exists a continuous path
\[
\gamma:[0,1]\to Y
\]
with
\[
\gamma(0)=y_0,
\qquad
\gamma(1)=y_1.
\]
Define
\[
H(x,t)=\gamma(t).
\]
Then
\[
H(x,0)=y_0,
\qquad
H(x,1)=y_1.
\]
Hence
\[
c_{y_0}\simeq c_{y_1}.
\]

\subsubsection{Example 5: The identity and antipodal map on \(S^1\)}

Regard
\[
S^1=\{z\in\mathbb C:|z|=1\}.
\]
Let
\[
f(z)=z,
\qquad
g(z)=-z.
\]
Define
\[
H(z,t)=e^{i\pi t}z.
\]
Then
\[
H(z,0)=z=f(z),
\qquad
H(z,1)=e^{i\pi}z=-z=g(z).
\]
Also \(|e^{i\pi t}z|=1\), so \(H(z,t)\in S^1\) for all \(z,t\). Therefore
\[
f\simeq g.
\]

\subsubsection{Example 6: A null-homotopic loop in \(\mathbb R^2\)}

Define
\[
f:[0,1]\to\mathbb R^2,
\qquad
f(s)=(\cos 2\pi s,\sin 2\pi s).
\]
This traces the unit circle in the plane. Since the target is \(\mathbb R^2\), the map is null-homotopic. Indeed, define
\[
H(s,t)=\big((1-t)\cos 2\pi s,\,(1-t)\sin 2\pi s\big).
\]
Then
\[
H(s,0)=f(s),
\qquad
H(s,1)=(0,0).
\]
Thus \(f\) is homotopic to the constant map at the origin.

\subsubsection{Example 7: A non-example on \(S^1\)}

Let
\[
f:S^1\to S^1,\qquad f(z)=z,
\]
and
\[
g:S^1\to S^1,\qquad g(z)=1.
\]
Then \(g\) is constant, but \(f\) is not homotopic to \(g\). Intuitively, \(f\) winds once around the circle, whereas \(g\) does not wind at all.

More formally,
\[
\deg(f)=1,
\qquad
\deg(g)=0.
\]
Since homotopic maps \(S^1\to S^1\) must have the same degree, \(f\not\simeq g\).

\subsubsection{Example 8: The antipodal map on \(S^n\)}

Let
\[
a:S^n\to S^n,\qquad a(x)=-x.
\]
Then:
\begin{itemize}
	\item if \(n\) is odd, \(a\simeq \mathrm{id}_{S^n}\),
	\item if \(n\) is even, \(a\not\simeq \mathrm{id}_{S^n}\).
\end{itemize}
This follows from degree theory, since
\[
\deg(a)=(-1)^{n+1}.
\]

\subsection{Basic properties}

Homotopy is an equivalence relation on the set of continuous maps \(X\to Y\):

\begin{enumerate}
	\item \textbf{Reflexive:} \(f\simeq f\), using the constant homotopy
	\[
	H(x,t)=f(x).
	\]
	
	\item \textbf{Symmetric:} if \(f\simeq g\), then \(g\simeq f\), by reversing time:
	\[
	\widetilde H(x,t)=H(x,1-t).
	\]
	
	\item \textbf{Transitive:} if \(f\simeq g\) and \(g\simeq h\), then \(f\simeq h\), by concatenating homotopies.
\end{enumerate}

\subsection{Homotopy equivalence of spaces}

Two spaces \(X\) and \(Y\) are called \emph{homotopy equivalent} if there exist continuous maps
\[
f:X\to Y,
\qquad
g:Y\to X
\]
such that
\[
g\circ f\simeq \mathrm{id}_X,
\qquad
f\circ g\simeq \mathrm{id}_Y.
\]

Examples include:
\begin{itemize}
	\item a disk and a point,
	\item an annulus and a circle,
	\item a circle with a line segment attached and a circle.
\end{itemize}

\section{Null-homotopic maps, homotopy equivalence, contractible spaces, and degree}

\subsection{Null-homotopic maps}

\subsubsection*{Definition}

Let \(X\) and \(Y\) be topological spaces, and let
\[
f:X\to Y
\]
be continuous. We say that \(f\) is \emph{null-homotopic} if \(f\) is homotopic to a constant map. That is, there exists a point \(y_0\in Y\) and a continuous map
\[
H:X\times[0,1]\to Y
\]
such that
\[
H(x,0)=f(x),\qquad H(x,1)=y_0
\]
for all \(x\in X\).

Equivalently,
\[
f\simeq c_{y_0},
\qquad\text{where }c_{y_0}(x)=y_0.
\]

\subsubsection*{Intuition}

A null-homotopy continuously deforms the whole image of the map into a single point. So a null-homotopic map is topologically inessential.

\subsubsection*{Examples}

\paragraph{Example 1.}
Every map
\[
f:X\to \mathbb R^n
\]
is null-homotopic, because
\[
H(x,t)=(1-t)f(x)
\]
deforms \(f\) to the constant map \(0\).

\paragraph{Example 2.}
The map
\[
f:[0,1]\to \mathbb R^2,\qquad f(s)=(\cos 2\pi s,\sin 2\pi s)
\]
is null-homotopic in \(\mathbb R^2\), via
\[
H(s,t)=\big((1-t)\cos 2\pi s,\,(1-t)\sin 2\pi s\big).
\]

\paragraph{Example 3.}
The same formula, regarded as a map into \(S^1\), is not null-homotopic, because it winds once around the circle.

\paragraph{Example 4.}
Every constant map is null-homotopic.

\subsection{Homotopy equivalence}

\subsubsection*{Definition}

Let \(X\) and \(Y\) be topological spaces. We say that \(X\) and \(Y\) are \emph{homotopy equivalent} if there exist continuous maps
\[
f:X\to Y,\qquad g:Y\to X
\]
such that
\[
g\circ f\simeq \mathrm{id}_X,
\qquad
f\circ g\simeq \mathrm{id}_Y.
\]

Then we write
\[
X\simeq Y.
\]

\subsubsection*{Intuition}

Homotopy equivalence means that two spaces have the same essential shape from the viewpoint of homotopy theory, even if they are not homeomorphic.

\subsubsection*{Examples}

\paragraph{Example 1.}
\(\mathbb R^n\) is homotopy equivalent to a point.

\paragraph{Example 2.}
An annulus
\[
A=\{(x,y)\in\mathbb R^2:1\le x^2+y^2\le 4\}
\]
is homotopy equivalent to \(S^1\), by radial deformation.

\paragraph{Example 3.}
A disk \(D^2\) is homotopy equivalent to a point.

\paragraph{Example 4.}
A circle with a line segment attached is homotopy equivalent to a circle.

\subsection{Contractible spaces}

\subsubsection*{Definition}

A space \(X\) is called \emph{contractible} if the identity map
\[
\mathrm{id}_X:X\to X
\]
is null-homotopic. Thus there exist \(x_0\in X\) and a continuous map
\[
H:X\times[0,1]\to X
\]
such that
\[
H(x,0)=x,\qquad H(x,1)=x_0
\]
for all \(x\in X\).

\subsubsection*{Intuition}

A contractible space can be continuously shrunk to a point inside itself.

\subsubsection*{Examples}

\paragraph{Example 1.}
\(\mathbb R^n\) is contractible via
\[
H(x,t)=(1-t)x.
\]

\paragraph{Example 2.}
The disk \(D^n\) is contractible via the same formula.

\paragraph{Example 3.}
Every convex subset of \(\mathbb R^n\) is contractible.

\paragraph{Example 4.}
Every star-shaped subset of \(\mathbb R^n\) is contractible.

\paragraph{Non-example 1.}
\(S^1\) is not contractible.

\paragraph{Non-example 2.}
\(S^n\) is not contractible for \(n\ge 1\).

\subsubsection*{Relation with homotopy equivalence}

A space \(X\) is contractible if and only if
\[
X\simeq \{*\}.
\]

\subsection{Degree of a map}

\subsubsection*{Definition}

Let
\[
f:S^n\to S^n
\]
be continuous. Since
\[
H_n(S^n)\cong \mathbb Z,
\]
the induced map
\[
f_*:H_n(S^n)\to H_n(S^n)
\]
must be multiplication by some integer. This integer is called the \emph{degree} of \(f\), and is denoted
\[
\deg(f).
\]

Thus, if \(\alpha\) is a generator of \(H_n(S^n)\), then
\[
f_*(\alpha)=\deg(f)\,\alpha.
\]

\subsubsection*{Intuition}

The degree measures how many times the sphere wraps around itself under the map, with sign recording orientation behavior.

\subsubsection*{Basic properties}

\begin{enumerate}
	\item If \(f\simeq g\), then
	\[
	\deg(f)=\deg(g).
	\]
	
	\item
	\[
	\deg(\mathrm{id}_{S^n})=1.
	\]
	
	\item If \(c:S^n\to S^n\) is constant, then
	\[
	\deg(c)=0.
	\]
	
	\item
	\[
	\deg(f\circ g)=\deg(f)\deg(g).
	\]
\end{enumerate}

\subsubsection*{Examples}

\paragraph{Example 1.}
The identity map on \(S^n\) has degree \(1\).

\paragraph{Example 2.}
A constant map has degree \(0\).

\paragraph{Example 3.}
The antipodal map
\[
a(x)=-x
\]
has degree
\[
\deg(a)=(-1)^{n+1}.
\]

\paragraph{Example 4.}
For
\[
f_m:S^1\to S^1,\qquad f_m(z)=z^m,
\]
we have
\[
\deg(f_m)=m.
\]

\paragraph{Example 5.}
The maps
\[
f(z)=z,\qquad g(z)=1
\]
from \(S^1\to S^1\) are not homotopic, because
\[
\deg(f)=1,\qquad \deg(g)=0.
\]

\subsection{Conclusion}

These four notions are closely related:

\begin{itemize}
	\item a map is \emph{null-homotopic} if it can be deformed to a constant map,
	\item two spaces are \emph{homotopy equivalent} if they have the same homotopy type,
	\item a space is \emph{contractible} if it is homotopy equivalent to a point,
	\item the \emph{degree} of a map \(S^n\to S^n\) is an integer homotopy invariant measuring global wrapping behavior.
\end{itemize}

\subsection{Conclusion}

A homotopy between maps \(f,g:X\to Y\) is a continuous deformation from \(f\) to \(g\):
\[
H:X\times[0,1]\to Y,
\qquad
H(x,0)=f(x),\quad H(x,1)=g(x).
\]
This notion is one of the central ideas of algebraic topology. It allows us to classify maps not by exact formulas, but by whether they can be continuously deformed into one another.	
	
	
\section{Problem 2 --- A map of \(S^1\) not homotopic to the identity}

\subsection{Exact statement}

Let
\[
f:S^1\to S^1
\]
be a map which is not homotopic to the identity map. Show that there exists an
\[
x\in S^1
\]
such that
\[
f(x)=x,
\]
and a
\[
y\in S^1
\]
such that
\[
f(y)=-y.
\]

\subsection{Solution}

We will prove separately that:

\begin{enumerate}
	\item there exists \(y\in S^1\) such that \(f(y)=-y\),
	\item there exists \(x\in S^1\) such that \(f(x)=x\).
\end{enumerate}

The key observation is that if \(f\) avoids antipodal points, then \(f\) is homotopic to the identity, and if \(f\) avoids fixed points, then \(f\) is homotopic to the antipodal map. Since on \(S^1\) the antipodal map is homotopic to the identity, either situation contradicts the assumption that \(f\) is not homotopic to the identity.

\subsubsection{Step 1: Existence of a point \(y\) such that \(f(y)=-y\)}

Assume, for contradiction, that
\[
f(z)\neq -z
\qquad\text{for all }z\in S^1.
\]

We identify \(S^1\) with the unit circle in \(\mathbb C\). Then for each \(z\in S^1\),
\[
\frac{f(z)}{z}\in S^1
\]
and the assumption \(f(z)\neq -z\) means precisely that
\[
\frac{f(z)}{z}\neq -1.
\]
Hence the map
\[
z\longmapsto \frac{f(z)}{z}
\]
takes values in
\[
S^1\setminus\{-1\},
\]
which is an open arc. Therefore there exists a continuous function
\[
\theta:S^1\to (-\pi,\pi)
\]
such that
\[
\frac{f(z)}{z}=e^{i\theta(z)}.
\]

Now define
\[
H:S^1\times[0,1]\to S^1,
\qquad
H(z,t)=e^{it\theta(z)}z.
\]

This map is continuous, and for all \(z\in S^1\),
\[
H(z,0)=z,
\]
while
\[
H(z,1)=e^{i\theta(z)}z=f(z).
\]
Thus \(H\) is a homotopy from \(\mathrm{id}_{S^1}\) to \(f\). Hence
\[
f\simeq \mathrm{id}_{S^1},
\]
contrary to the assumption.

Therefore our assumption was false. So there exists \(y\in S^1\) such that
\[
f(y)=-y.
\]

\subsubsection{Step 2: Existence of a point \(x\) such that \(f(x)=x\)}

Assume, for contradiction, that
\[
f(z)\neq z
\qquad\text{for all }z\in S^1.
\]

Let
\[
a:S^1\to S^1,\qquad a(z)=-z
\]
be the antipodal map. We claim that \(f\) is homotopic to \(a\).

Indeed, for each \(z\in S^1\),
\[
\frac{f(z)}{-z}\in S^1.
\]
The assumption \(f(z)\neq z\) implies
\[
\frac{f(z)}{-z}\neq -1,
\]
because
\[
\frac{f(z)}{-z}=-1
\quad\Longleftrightarrow\quad
f(z)=z.
\]
Thus the map
\[
z\longmapsto \frac{f(z)}{-z}
\]
also takes values in \(S^1\setminus\{-1\}\). Therefore there exists a continuous function
\[
\phi:S^1\to (-\pi,\pi)
\]
such that
\[
\frac{f(z)}{-z}=e^{i\phi(z)}.
\]

Define
\[
K:S^1\times[0,1]\to S^1,
\qquad
K(z,t)=e^{it\phi(z)}(-z).
\]

Then
\[
K(z,0)=-z=a(z),
\]
and
\[
K(z,1)=e^{i\phi(z)}(-z)=f(z).
\]
So \(K\) is a homotopy from \(a\) to \(f\). Therefore
\[
f\simeq a.
\]

But on \(S^1\), the antipodal map is homotopic to the identity. For example,
\[
L(z,t)=e^{i\pi t}z
\]
defines a homotopy from \(z\) to \(-z\). Hence
\[
a\simeq \mathrm{id}_{S^1}.
\]
It follows that
\[
f\simeq a\simeq \mathrm{id}_{S^1},
\]
again contradicting the assumption.

Therefore our assumption was false. So there exists \(x\in S^1\) such that
\[
f(x)=x.
\]

\subsection{Conclusion}

We have proved that if
\[
f:S^1\to S^1
\]
is not homotopic to the identity, then:

\begin{itemize}
	\item there exists \(x\in S^1\) such that
	\[
	f(x)=x,
	\]
	\item there exists \(y\in S^1\) such that
	\[
	f(y)=-y.
	\]
\end{itemize}

This is exactly the required statement.

\section{Problem 3 --- From one-sided homotopy inverses to a homotopy equivalence}

\subsection{Exact statement}

Suppose
\[
f:X\to Y
\]
is a map for which there exist maps
\[
g,h:Y\to X
\]
such that
\[
g\circ f \simeq \operatorname{Id}_X
\qquad\text{and}\qquad
f\circ h \simeq \operatorname{Id}_Y.
\]
Show that \(f\), \(g\), and \(h\) are homotopy equivalences.

\subsection{Idea}

We are given a left homotopy inverse \(g\) of \(f\) and a right homotopy inverse \(h\) of \(f\).  
The key point is to show that
\[
g\simeq h.
\]
Then \(g\) becomes both a left and a right homotopy inverse of \(f\), so \(f\) is a homotopy equivalence.

\subsection{Proof}

We have
\[
h \simeq \operatorname{Id}_X\circ h.
\]
Since
\[
g\circ f \simeq \operatorname{Id}_X,
\]
it follows that
\[
\operatorname{Id}_X\circ h \simeq (g\circ f)\circ h.
\]
By associativity,
\[
(g\circ f)\circ h = g\circ(f\circ h).
\]
Since
\[
f\circ h \simeq \operatorname{Id}_Y,
\]
we obtain
\[
g\circ(f\circ h)\simeq g\circ\operatorname{Id}_Y = g.
\]
Therefore
\[
h\simeq g.
\]

Now we already know
\[
g\circ f \simeq \operatorname{Id}_X.
\]
Also,
\[
f\circ g \simeq f\circ h \simeq \operatorname{Id}_Y.
\]
Hence \(g\) is a homotopy inverse of \(f\), so \(f\) is a homotopy equivalence.

Then \(g\) is also a homotopy equivalence, with homotopy inverse \(f\).

Finally, since \(h\simeq g\), we get
\[
h\circ f \simeq g\circ f \simeq \operatorname{Id}_X
\]
and by assumption
\[
f\circ h \simeq \operatorname{Id}_Y.
\]
So \(h\) is also a homotopy inverse of \(f\), and thus \(h\) is a homotopy equivalence.

Therefore \(f\), \(g\), and \(h\) are all homotopy equivalences.

\subsection{Interpretation}

This shows that if a map has a left inverse up to homotopy and a right inverse up to homotopy, then these inverses agree up to homotopy, and the map is a genuine homotopy equivalence.

\section{Problem 4 --- A retract of a contractible space is contractible}

\subsection{Exact statement}

Show that a retract of a contractible space is contractible.

\subsection{Definitions}

A space \(A\) is called a retract of a space \(X\) if there exist maps
\[
i:A\to X,
\qquad
r:X\to A
\]
such that
\[
r\circ i = \operatorname{Id}_A.
\]

A space \(X\) is contractible if
\[
\operatorname{Id}_X \simeq c
\]
for some constant map \(c:X\to X\).

\subsection{Idea}

If \(X\) contracts to a point, then a retract \(A\) of \(X\) also contracts: perform the contraction in \(X\), then apply the retraction back to \(A\).

\subsection{Proof}

Let \(A\) be a retract of \(X\). Then there exist maps
\[
i:A\to X,
\qquad
r:X\to A
\]
such that
\[
r\circ i = \operatorname{Id}_A.
\]

Assume that \(X\) is contractible. Then there exists a point \(x_0\in X\) such that
\[
\operatorname{Id}_X \simeq c_{x_0},
\]
where \(c_{x_0}:X\to X\) is the constant map with value \(x_0\).

Compose this homotopy on the left with \(r\) and on the right with \(i\). Since homotopy is preserved under composition, we get
\[
r\circ \operatorname{Id}_X \circ i \simeq r\circ c_{x_0}\circ i.
\]
But
\[
r\circ \operatorname{Id}_X \circ i = r\circ i = \operatorname{Id}_A.
\]
Hence
\[
\operatorname{Id}_A \simeq r\circ c_{x_0}\circ i.
\]
Since \(c_{x_0}\) is constant, the composite \(r\circ c_{x_0}\circ i\) is also constant, with value \(r(x_0)\in A\).

Therefore \(\operatorname{Id}_A\) is homotopic to a constant map, so \(A\) is contractible.

\subsection{Remark}

This result is often used as an obstruction: if a space \(A\) is not contractible, then it cannot be a retract of any contractible space.

\section{Examples and Counterexamples for Problems 3 and 4}

\subsection{Problem 3 --- One-sided homotopy inverses imply homotopy equivalence}

\subsubsection{Statement recalled}

Suppose
\[
f:X\to Y
\]
is a map for which there exist maps
\[
g,h:Y\to X
\]
such that
\[
g\circ f \simeq \operatorname{id}_X
\qquad\text{and}\qquad
f\circ h \simeq \operatorname{id}_Y.
\]
Then \(f\), \(g\), and \(h\) are homotopy equivalences.

\subsubsection{Why examples matter}

This theorem says that having a left inverse up to homotopy and a right inverse up to homotopy is enough to force a genuine homotopy equivalence.

The most useful way to understand it is:

\begin{itemize}
	\item first, through geometric examples where the theorem applies;
	\item second, through counterexamples showing that only one of the two assumptions is not enough.
\end{itemize}

\subsubsection{Example 1: the cylinder projects onto the circle}

Let
\[
X=S^1\times I,
\qquad
Y=S^1,
\]
where \(I=[0,1]\). Define
\[
f:S^1\times I\to S^1,
\qquad
f(z,t)=z.
\]
Define also
\[
g:S^1\to S^1\times I,
\qquad
g(z)=(z,0),
\]
and
\[
h:S^1\to S^1\times I,
\qquad
h(z)=(z,1).
\]

Then
\[
f\circ h(z)=f(z,1)=z=\operatorname{id}_{S^1}(z),
\]
so
\[
f\circ h=\operatorname{id}_{S^1}.
\]

On the other hand,
\[
g\circ f(z,t)=g(z)=(z,0).
\]
This is not equal to \(\operatorname{id}_{S^1\times I}\), but it is homotopic to it. Indeed, define
\[
H:(S^1\times I)\times I\to S^1\times I,
\qquad
H((z,t),s)=(z,(1-s)t).
\]
Then
\[
H((z,t),0)=(z,t)=\operatorname{id}_{S^1\times I}(z,t),
\]
and
\[
H((z,t),1)=(z,0)=g\circ f(z,t).
\]
Hence
\[
g\circ f \simeq \operatorname{id}_{S^1\times I}.
\]

So the hypotheses of the theorem hold. Therefore \(f\) is a homotopy equivalence, and both \(g\) and \(h\) are homotopy inverses of \(f\) up to homotopy.

Notice also that \(g\) and \(h\) are themselves homotopic, via
\[
K:S^1\times I\to S^1\times I,
\qquad
K(z,s)=(z,s),
\]
viewed as a homotopy from \(g\) to \(h\), since
\[
K(z,0)=g(z),
\qquad
K(z,1)=h(z).
\]

\bigskip

\textbf{Geometric meaning.}
The cylinder \(S^1\times I\) has the same homotopy type as the circle \(S^1\). The projection forgets the interval direction, and the inclusions of the bottom and top circles both give homotopy inverses.

\subsubsection{Example 2: deformation retractions}

A very important source of examples comes from deformation retracts.

Suppose \(Y\subseteq X\), and suppose there exists a retraction
\[
r:X\to Y
\]
such that the composite
\[
i\circ r:X\to X
\]
is homotopic to \(\operatorname{id}_X\), where
\[
i:Y\hookrightarrow X
\]
is the inclusion. Then \(Y\) is a deformation retract of \(X\).

Now set
\[
f=r:X\to Y,
\qquad
g=i:Y\to X,
\qquad
h=i:Y\to X.
\]
Then
\[
f\circ h=r\circ i=\operatorname{id}_Y,
\]
and
\[
g\circ f=i\circ r \simeq \operatorname{id}_X.
\]
So the theorem applies immediately.

Thus every deformation retraction produces an example of the theorem.

\paragraph{Concrete instance.}
Take
\[
X=\mathbb{R}^2\setminus\{0\},
\qquad
Y=S^1.
\]
Define
\[
f(x)=\frac{x}{\|x\|},
\qquad
g(y)=y.
\]
Then
\[
f\circ g(y)=\frac{y}{\|y\|}=y,
\]
because \(y\in S^1\) has norm \(1\). So
\[
f\circ g=\operatorname{id}_{S^1}.
\]
Also
\[
g\circ f(x)=\frac{x}{\|x\|}.
\]
This is homotopic to \(\operatorname{id}_{\mathbb{R}^2\setminus\{0\}}\) by the homotopy
\[
H(x,s)=\left((1-s)+\frac{s}{\|x\|}\right)x.
\]
Indeed,
\[
H(x,0)=x,
\qquad
H(x,1)=\frac{x}{\|x\|}.
\]
Since \(x\neq 0\), the coefficient
\[
(1-s)+\frac{s}{\|x\|}
\]
is always positive, so \(H(x,s)\neq 0\) for all \(s\), and the homotopy stays inside \(\mathbb{R}^2\setminus\{0\}\).

Hence \(S^1\) and \(\mathbb{R}^2\setminus\{0\}\) are homotopy equivalent.

\subsubsection{Example 3: a contractible space and a point}

Let
\[
X=I=[0,1],
\qquad
Y=\{*\},
\]
and let
\[
f:X\to Y
\]
be the unique map.

Define
\[
g(*)=0,
\qquad
h(*)=1.
\]
Then
\[
f\circ h=\operatorname{id}_{\{*\}},
\]
because any map from a one-point space to itself is the identity.

Also
\[
g\circ f:X\to X
\]
is the constant map at \(0\). Since the interval is contractible, this constant map is homotopic to \(\operatorname{id}_I\). For example, the homotopy
\[
H:I\times I\to I,
\qquad
H(x,s)=(1-s)x
\]
satisfies
\[
H(x,0)=x,
\qquad
H(x,1)=0.
\]
Thus
\[
g\circ f \simeq \operatorname{id}_I.
\]

So the theorem applies. In particular, the unique map from a contractible space to a point is a homotopy equivalence.

This is a very instructive example because \(g\) and \(h\) are genuinely different maps, but they are still homotopic.

\subsubsection{Counterexample 1: a left homotopy inverse alone is not enough}

We now show that the condition
\[
g\circ f \simeq \operatorname{id}_X
\]
by itself does \emph{not} imply that \(f\) is a homotopy equivalence.

Take
\[
X=S^1,
\qquad
Y=S^1\vee S^2,
\]
the wedge of a circle and a \(2\)-sphere. Let
\[
f:S^1\hookrightarrow S^1\vee S^2
\]
be the inclusion of the \(S^1\)-summand.

Define
\[
g:S^1\vee S^2\to S^1
\]
to be the identity on the \(S^1\)-summand and constant on the \(S^2\)-summand, sending the sphere to the wedge point.

Then
\[
g\circ f=\operatorname{id}_{S^1}.
\]
So \(f\) has a left inverse, even exactly, not only up to homotopy.

However, \(f\) is \emph{not} a homotopy equivalence. Indeed,
\[
H_2(S^1)=0,
\qquad
H_2(S^1\vee S^2)\cong \mathbb{Z}.
\]
Homotopy equivalent spaces must have isomorphic homology groups, but these do not. Therefore \(S^1\) and \(S^1\vee S^2\) are not homotopy equivalent.

So a left inverse up to homotopy alone is not enough.

\subsubsection{Counterexample 2: a right homotopy inverse alone is not enough}

Now we show that the condition
\[
f\circ h \simeq \operatorname{id}_Y
\]
alone is also not enough.

Take
\[
X=S^1,
\qquad
Y=\{*\},
\]
and let
\[
f:S^1\to \{*\}
\]
be the unique map. Let
\[
h:\{*\}\to S^1
\]
choose any point of \(S^1\).

Then
\[
f\circ h=\operatorname{id}_{\{*\}},
\]
again because any map from a one-point space to itself is the identity.

But \(f\) is not a homotopy equivalence, because that would imply that \(S^1\) is homotopy equivalent to a point, hence contractible, which is false.

So a right inverse up to homotopy alone is not enough either.

\subsubsection{Conclusion for Problem 3}

The theorem is strong precisely because it uses \emph{both} assumptions:
\[
g\circ f \simeq \operatorname{id}_X
\qquad\text{and}\qquad
f\circ h \simeq \operatorname{id}_Y.
\]
Neither one alone suffices, but together they force \(f\) to be a homotopy equivalence.

\bigskip

\textbf{Solution.}
The examples above show the theorem in natural geometric settings: cylinders, deformation retracts, and contractible spaces. The counterexamples show that the existence of only a left inverse or only a right inverse up to homotopy is too weak to guarantee homotopy equivalence.

\subsection{Problem 4 --- A retract of a contractible space is contractible}

\subsubsection{Statement recalled}

A retract of a contractible space is contractible.

That is, if \(A\) is a retract of \(X\), meaning there exist maps
\[
i:A\to X,
\qquad
r:X\to A
\]
such that
\[
r\circ i=\operatorname{id}_A,
\]
and if \(X\) is contractible, then \(A\) is contractible.

\subsubsection{Why examples matter}

This theorem tells us that retracts inherit contractibility. It is one of the simplest and most useful tools for proving that certain subspaces cannot be retracts.

To understand it well, one should compare:

\begin{itemize}
	\item examples where the retract really is contractible;
	\item counterexamples showing that simply being a subspace of a contractible space is not enough.
\end{itemize}

\subsubsection{Example 1: a point is a retract of an interval}

Let
\[
A=\{0\},
\qquad
X=[0,1].
\]
Define the inclusion
\[
i:A\hookrightarrow X,
\]
and define
\[
r:X\to A,
\qquad
r(x)=0.
\]
Then
\[
r\circ i(0)=r(0)=0,
\]
so
\[
r\circ i=\operatorname{id}_A.
\]
Hence \(A\) is a retract of \([0,1]\).

Since \([0,1]\) is contractible, the theorem tells us that \(A\) is contractible.

Of course a point is obviously contractible, but this is the simplest illustration of the theorem.

\subsubsection{Example 2: a diameter is a retract of the disk}

Let
\[
X=D^2=\{(x,y)\in\mathbb{R}^2:x^2+y^2\le 1\},
\]
and let
\[
A=[-1,1]\times\{0\}\subseteq D^2.
\]
Define
\[
r:D^2\to A,
\qquad
r(x,y)=(x,0).
\]
This is continuous. Moreover, if \((x,0)\in A\), then
\[
r(x,0)=(x,0),
\]
so \(r\) restricts to the identity on \(A\). Thus
\[
r|_A=\operatorname{id}_A,
\]
and equivalently
\[
r\circ i=\operatorname{id}_A,
\]
where \(i:A\hookrightarrow D^2\) is the inclusion.

Therefore \(A\) is a retract of the disk. Since the disk is contractible, the theorem implies that \(A\) is contractible.

This gives a geometric picture: collapsing vertical directions in the disk gives a retraction onto a line segment.

\subsubsection{Example 3: \(\mathbb{R}^k\) is a retract of \(\mathbb{R}^n\)}

Assume \(k\le n\). Identify
\[
\mathbb{R}^k
\cong
\mathbb{R}^k\times\{0\}\subseteq \mathbb{R}^n.
\]
Define
\[
r:\mathbb{R}^n\to \mathbb{R}^k\times\{0\},
\qquad
r(x_1,\dots,x_n)=(x_1,\dots,x_k,0,\dots,0).
\]
This is continuous and fixes every point of \(\mathbb{R}^k\times\{0\}\). So it is a retraction.

Since \(\mathbb{R}^n\) is contractible, the theorem implies that \(\mathbb{R}^k\) is contractible.

Again, this is obvious independently, but it is a clean example of the theorem in linear geometry.

\subsubsection{Example 4: any retract of a convex set is contractible}

Let \(C\subseteq \mathbb{R}^n\) be convex. Then \(C\) is contractible, because for any chosen point \(p\in C\), the homotopy
\[
H:C\times I\to C,
\qquad
H(x,s)=(1-s)x+sp
\]
contracts \(C\) to \(p\).

Therefore, if \(A\subseteq C\) is a retract of \(C\), then \(A\) is contractible.

This gives a broad family of examples. In practice, whenever one has a retraction from a convex set onto a subspace, that subspace must be contractible.

\subsubsection{Counterexample 1: a subspace of a contractible space need not be contractible}

The theorem does \emph{not} say that every subspace of a contractible space is contractible.

A standard counterexample is
\[
S^1\subseteq D^2.
\]
The disk \(D^2\) is contractible, but \(S^1\) is not contractible.

Thus the statement
\[
\text{``every subspace of a contractible space is contractible''}
\]
is false.

This shows why the word \emph{retract} is essential.

\subsubsection{Counterexample 2: \(S^1\) is not a retract of \(D^2\)}

The previous counterexample can be sharpened.

Since \(D^2\) is contractible and \(S^1\) is not contractible, the theorem implies that \(S^1\) cannot be a retract of \(D^2\).

So although
\[
S^1\subseteq D^2,
\]
there is no continuous retraction
\[
r:D^2\to S^1
\]
such that \(r|_{S^1}=\operatorname{id}_{S^1}\).

This is one of the classical uses of the theorem: it gives a quick proof that the boundary circle is not a retract of the disk.

\subsubsection{Counterexample 3: a continuous image of a contractible space need not be contractible}

The theorem also does not say that every continuous image of a contractible space is contractible.

Consider the map
\[
f:[0,1]\to S^1,
\qquad
f(t)=e^{2\pi i t}.
\]
The interval \([0,1]\) is contractible, and \(f\) is surjective, so \(S^1\) is a continuous image of a contractible space.

But \(S^1\) is not contractible.

Therefore the statement
\[
\text{``every continuous image of a contractible space is contractible''}
\]
is false.

This is another way to see that ``retract'' is much stronger than merely ``image''.

\subsubsection{Counterexample 4: a homotopy equivalent copy is different from a retract}

It is also useful to distinguish retracts from homotopy equivalent subspaces.

For example, the circle \(S^1\) is homotopy equivalent to the annulus
\[
A=\{x\in\mathbb{R}^2:1\le \|x\|\le 2\},
\]
because the annulus deformation retracts onto its middle circle.

But \(S^1\) is not a retract of the disk \(D^2\), even though both \(S^1\) and the annulus appear naturally inside the plane.

So one must not confuse:

\begin{itemize}
	\item being homotopy equivalent to something contractible or noncontractible,
	\item being contained in a contractible space,
	\item being a retract of a contractible space.
\end{itemize}

The retract condition is much more rigid.

\subsubsection{Conclusion for Problem 4}

The theorem says:
\[
\text{retract of a contractible space } \Longrightarrow \text{ contractible}.
\]

But the following weaker statements are false:

\begin{itemize}
	\item every subspace of a contractible space is contractible;
	\item every continuous image of a contractible space is contractible.
\end{itemize}

So the hypothesis that we have a genuine retraction is exactly what makes the theorem true.

\bigskip

\textbf{Solution.}
The examples above illustrate the theorem in simple geometric settings such as intervals, disks, Euclidean spaces, and convex sets. The counterexamples show that the condition ``retract'' cannot be weakened to ``subspace'' or ``continuous image.'' In particular, the noncontractibility of \(S^1\) proves immediately that \(S^1\) cannot be a retract of \(D^2\).

\subsection{Final summary}

For Problem 3, the theorem becomes much clearer when one compares:

\begin{itemize}
	\item positive examples coming from cylinders, deformation retracts, and contractible spaces;
	\item negative examples showing that a left inverse alone or a right inverse alone is insufficient.
\end{itemize}

For Problem 4, the theorem is best understood by comparing:

\begin{itemize}
	\item retracts of contractible spaces, which must themselves be contractible;
	\item arbitrary subspaces or images of contractible spaces, which need not be contractible.
\end{itemize}

These examples and counterexamples are often more instructive than the abstract statement itself, because they show exactly where the hypotheses are used and why they are necessary.

\section{Definitions Used in Problems 3 and 4}

\subsection{Homotopies, homotopy equivalences, retracts, and contractions}

In Problems \(3\) and \(4\), several standard notions from homotopy theory were used.  
We collect them here in one place, together with the concrete constructions that appeared in the examples.

\subsubsection{Homotopy of maps}

Let \(X\) and \(Y\) be topological spaces, and let
\[
f,g:X\to Y
\]
be continuous maps.

We say that \(f\) is \emph{homotopic} to \(g\), written
\[
f\simeq g,
\]
if there exists a continuous map
\[
H:X\times I\to Y,
\qquad I=[0,1],
\]
such that
\[
H(x,0)=f(x),
\qquad
H(x,1)=g(x)
\]
for all \(x\in X\).

Such a map \(H\) is called a \emph{homotopy} from \(f\) to \(g\).

\paragraph{Interpretation.}
A homotopy is a continuous deformation of one map into another, with the parameter \(s\in I\) representing time.

\paragraph{Example.}
On the interval \(I=[0,1]\), the identity map and the constant map at \(0\) are homotopic via
\[
H:I\times I\to I,
\qquad
H(x,s)=(1-s)x.
\]
Indeed,
\[
H(x,0)=x=\operatorname{id}_I(x),
\qquad
H(x,1)=0.
\]

\subsubsection{Identity map}

For any space \(X\), the \emph{identity map}
\[
\operatorname{id}_X:X\to X
\]
is defined by
\[
\operatorname{id}_X(x)=x
\qquad
\text{for all }x\in X.
\]

This is the map that leaves every point fixed.

\subsubsection{Constant map}

Let \(X\) and \(Y\) be spaces, and let \(y_0\in Y\).  
The \emph{constant map} with value \(y_0\) is the map
\[
c_{y_0}:X\to Y
\]
defined by
\[
c_{y_0}(x)=y_0
\qquad
\text{for all }x\in X.
\]

So a constant map sends every point of \(X\) to the same single point of \(Y\).

\paragraph{Example.}
If \(X=[0,1]\) and \(Y=[0,1]\), then the map
\[
c_0(x)=0
\]
for all \(x\in[0,1]\) is the constant map at \(0\).

\subsubsection{Homotopy equivalence}

A map
\[
f:X\to Y
\]
is called a \emph{homotopy equivalence} if there exists a continuous map
\[
g:Y\to X
\]
such that
\[
g\circ f \simeq \operatorname{id}_X
\qquad\text{and}\qquad
f\circ g \simeq \operatorname{id}_Y.
\]

The map \(g\) is then called a \emph{homotopy inverse} of \(f\).

If such maps exist, we say that \(X\) and \(Y\) are \emph{homotopy equivalent}, and write
\[
X\simeq Y.
\]

\paragraph{Interpretation.}
Homotopy equivalent spaces may not be homeomorphic, but they have the same shape from the viewpoint of homotopy theory.

\paragraph{Example.}
The cylinder \(S^1\times I\) is homotopy equivalent to the circle \(S^1\).  
The projection
\[
f:S^1\times I\to S^1,
\qquad
f(z,t)=z,
\]
is a homotopy equivalence.

\subsubsection{Left and right homotopy inverses}

Let
\[
f:X\to Y.
\]

A map
\[
g:Y\to X
\]
is called a \emph{left homotopy inverse} of \(f\) if
\[
g\circ f \simeq \operatorname{id}_X.
\]

A map
\[
h:Y\to X
\]
is called a \emph{right homotopy inverse} of \(f\) if
\[
f\circ h \simeq \operatorname{id}_Y.
\]

Problem \(3\) showed that if both a left homotopy inverse and a right homotopy inverse exist, then \(f\) is actually a homotopy equivalence.

\subsubsection{Retract}

Let \(A\subseteq X\). We say that \(A\) is a \emph{retract} of \(X\) if there exists a continuous map
\[
r:X\to A
\]
such that
\[
r(a)=a
\qquad
\text{for all }a\in A.
\]

Equivalently, if
\[
i:A\hookrightarrow X
\]
denotes the inclusion, then
\[
r\circ i=\operatorname{id}_A.
\]

The map \(r\) is called a \emph{retraction}.

\paragraph{Interpretation.}
A retract is a subspace onto which the whole space can be continuously projected while keeping that subspace pointwise fixed.

\paragraph{Example.}
Let
\[
A=[-1,1]\times\{0\}\subseteq D^2.
\]
Then
\[
r:D^2\to A,
\qquad
r(x,y)=(x,0),
\]
is a retraction. Hence \(A\) is a retract of the disk.

\subsubsection{Deformation retract}

Let \(A\subseteq X\). We say that \(A\) is a \emph{deformation retract} of \(X\) if there exists a retraction
\[
r:X\to A
\]
such that the composite
\[
i\circ r:X\to X
\]
is homotopic to \(\operatorname{id}_X\), where \(i:A\hookrightarrow X\) is the inclusion.

Equivalently, there exists a homotopy
\[
H:X\times I\to X
\]
such that
\[
H(x,0)=x,
\qquad
H(x,1)\in A,
\qquad
H(a,1)=a
\quad\text{for all }a\in A.
\]

\paragraph{Interpretation.}
A deformation retract is stronger than a retract: the whole space can be continuously shrunk onto the subspace.

\paragraph{Example.}
The circle \(S^1\) is a deformation retract of \(\mathbb{R}^2\setminus\{0\}\) via the map
\[
r(x)=\frac{x}{\|x\|}.
\]

\subsubsection{Contractible space}

A space \(X\) is called \emph{contractible} if the identity map is homotopic to a constant map. That is, there exists a point \(x_0\in X\) such that
\[
\operatorname{id}_X \simeq c_{x_0}.
\]

Equivalently, there exists a homotopy
\[
H:X\times I\to X
\]
such that
\[
H(x,0)=x,
\qquad
H(x,1)=x_0
\]
for all \(x\in X\).

\paragraph{Interpretation.}
A contractible space can be continuously shrunk to a single point.

\paragraph{Example.}
The interval \(I=[0,1]\) is contractible.  
A contraction to \(0\) is given by
\[
H:I\times I\to I,
\qquad
H(x,s)=(1-s)x.
\]

\subsubsection{Contraction}

A \emph{contraction} of a space \(X\) to a point \(x_0\in X\) is precisely a homotopy
\[
H:X\times I\to X
\]
from \(\operatorname{id}_X\) to the constant map \(c_{x_0}\). Thus
\[
H(x,0)=x,
\qquad
H(x,1)=x_0
\]
for all \(x\in X\).

So to say that a space is contractible is exactly to say that it admits a contraction.

\subsubsection{Contractions used in the examples}

We now record explicitly the contractions and homotopies used earlier.

\paragraph{1. Contraction of the interval.}
For \(X=I=[0,1]\), a contraction to \(0\) is
\[
H(x,s)=(1-s)x.
\]
At time \(s=0\), this is the identity:
\[
H(x,0)=x.
\]
At time \(s=1\), every point has moved to \(0\):
\[
H(x,1)=0.
\]

\paragraph{2. Contraction of a convex subset of \(\mathbb{R}^n\).}
Let \(C\subseteq \mathbb{R}^n\) be convex, and choose \(p\in C\). Define
\[
H:C\times I\to C,
\qquad
H(x,s)=(1-s)x+sp.
\]
Because \(C\) is convex, every point of the line segment joining \(x\) and \(p\) lies in \(C\), so \(H\) is well-defined. Also,
\[
H(x,0)=x,
\qquad
H(x,1)=p.
\]
Thus \(C\) is contractible.

\paragraph{3. Homotopy on the cylinder.}
For
\[
X=S^1\times I,
\]
we used the homotopy
\[
H((z,t),s)=(z,(1-s)t)
\]
from the identity map to the map \((z,t)\mapsto (z,0)\). Indeed,
\[
H((z,t),0)=(z,t),
\qquad
H((z,t),1)=(z,0).
\]

This is not a contraction of the cylinder to a point, but it is a deformation of the cylinder onto its bottom circle.

\paragraph{4. Radial deformation of \(\mathbb{R}^2\setminus\{0\}\) onto \(S^1\).}
For
\[
X=\mathbb{R}^2\setminus\{0\},
\]
we used the homotopy
\[
H(x,s)=\left((1-s)+\frac{s}{\|x\|}\right)x
\]
from the identity map to the radial projection
\[
x\mapsto \frac{x}{\|x\|}.
\]
Indeed,
\[
H(x,0)=x,
\qquad
H(x,1)=\frac{x}{\|x\|}.
\]

This shows that \(\mathbb{R}^2\setminus\{0\}\) deformation retracts onto \(S^1\).

\subsubsection{How these definitions fit Problems 3 and 4}

Problem \(3\) used the notions of:

\begin{itemize}
	\item homotopy of maps,
	\item identity map,
	\item constant map only indirectly,
	\item left and right homotopy inverses,
	\item homotopy equivalence.
\end{itemize}

Problem \(4\) used the notions of:

\begin{itemize}
	\item retract,
	\item retraction,
	\item contractible space,
	\item contraction.
\end{itemize}

The key logical ideas were:

\begin{itemize}
	\item in Problem \(3\), a left homotopy inverse and a right homotopy inverse together force a homotopy equivalence;
	\item in Problem \(4\), a contraction of the big space can be composed with a retraction to obtain a contraction of the retract.
\end{itemize}

\subsubsection{Final remark}

These definitions are basic, but they are extremely important.  
Many theorems in algebraic topology are really statements about how these notions interact:

\begin{itemize}
	\item homotopies compare maps,
	\item homotopy equivalences compare spaces,
	\item retracts identify rigidly embedded subspaces,
	\item contractions detect spaces with trivial homotopy type.
\end{itemize}

Once these definitions are fully understood, Problems \(3\) and \(4\) become natural and almost inevitable.
	

\section{Problem 5 - Strong Deformation Retractions and a Space Containing Both an Annulus and a Möbius Band as Deformation Retracts}

\textbf{Statement.}
Show that if a space \(X\) strongly deformation retracts to a point \(x_0 \in X\), then for every open neighbourhood \(x_0 \in U\) there exists a smaller open neighbourhood \(x_0 \in V \subset U\) such that the inclusion
\[
(V,x_0)\hookrightarrow (U,x_0)
\]
is based homotopic to the constant map.

\medskip

\textbf{Intuition.}
A strong deformation retraction gives a homotopy
\[
H:X\times I\to X
\]
from the identity to the constant map at \(x_0\), and throughout the homotopy the point \(x_0\) stays fixed:
\[
H(x_0,t)=x_0.
\]
Since \(x_0\) stays inside \(U\) for all \(t\), continuity implies that points sufficiently close to \(x_0\) also stay inside \(U\) for all \(t\). That smaller neighbourhood is the required \(V\).

\medskip

\textbf{Solution.}
Let
\[
H:X\times I\to X
\]
be a strong deformation retraction of \(X\) onto the point \(x_0\). Thus
\[
H(x,0)=x,
\qquad
H(x,1)=x_0,
\qquad
H(x_0,t)=x_0
\]
for all \(x\in X\) and all \(t\in I\).

Now let \(U\) be any open neighbourhood of \(x_0\).

For each \(t\in I\), we have
\[
H(x_0,t)=x_0\in U.
\]
Since \(H\) is continuous at the point \((x_0,t)\), there exist an open neighbourhood \(W_t\) of \(x_0\) in \(X\) and an open interval \(J_t\subset I\) containing \(t\) such that
\[
H(W_t\times J_t)\subset U.
\]

The family \(\{J_t\}_{t\in I}\) is an open cover of the compact interval \(I=[0,1]\). Hence there exist finitely many points
\[
t_1,\dots,t_n\in I
\]
such that
\[
I=J_{t_1}\cup\cdots\cup J_{t_n}.
\]

Define
\[
V=\Bigl(\bigcap_{i=1}^n W_{t_i}\Bigr)\cap U.
\]
Then \(V\) is an open neighbourhood of \(x_0\), and clearly
\[
V\subset U.
\]

We now claim that
\[
H(V\times I)\subset U.
\]
Indeed, let \(v\in V\) and \(s\in I\). Since the intervals \(J_{t_1},\dots,J_{t_n}\) cover \(I\), there exists some index \(i\) such that
\[
s\in J_{t_i}.
\]
Also,
\[
v\in V\subset W_{t_i}.
\]
Therefore
\[
H(v,s)\in H(W_{t_i}\times J_{t_i})\subset U.
\]
This proves the claim.

Hence the restriction
\[
K=H|_{V\times I}:V\times I\to U
\]
is well-defined. Moreover, for all \(v\in V\) and \(t\in I\),
\[
K(v,0)=H(v,0)=v,
\]
\[
K(v,1)=H(v,1)=x_0,
\]
and
\[
K(x_0,t)=H(x_0,t)=x_0.
\]
Thus \(K\) is a based homotopy in \((U,x_0)\) from the inclusion
\[
(V,x_0)\hookrightarrow (U,x_0)
\]
to the constant map at \(x_0\).

Therefore the inclusion
\[
(V,x_0)\hookrightarrow (U,x_0)
\]
is based homotopic to the constant map.

\section{Problem 6}

\textbf{Statement.}
Construct a space which contains both the annulus \(S^1\times I\) and the Möbius band as deformation retracts.

\medskip

\textbf{Idea.}
Both the annulus and the Möbius band deformation retract onto a circle, namely their core circle. So we glue an annulus and a Möbius band along their core circles. Then collapsing one of them onto the common core circle while keeping the other fixed gives a deformation retraction onto the other.

\medskip

\textbf{Construction of the space.}
Let \(A\) be an annulus and \(M\) a Möbius band. Let
\[
C_A\subset A
\qquad\text{and}\qquad
C_M\subset M
\]
be their core circles.

Choose a homeomorphism
\[
\varphi:C_A\to C_M.
\]
Now form the quotient space
\[
Y=(A\sqcup M)/\sim,
\]
where \(c\sim \varphi(c)\) for all \(c\in C_A\).

Equivalently, we identify the core circle of the annulus with the core circle of the Möbius band. Denote the common image of these circles by
\[
C\subset Y.
\]
Then \(Y\) contains a copy of the annulus \(A\) and a copy of the Möbius band \(M\), intersecting exactly in the common circle \(C\).

\medskip

\textbf{Why \(Y\) deformation retracts onto the annulus.}
The Möbius band deformation retracts onto its core circle. Let
\[
r_M:M\times I\to M
\]
be such a deformation retraction. Thus
\[
r_M(m,0)=m,
\qquad
r_M(m,1)\in C,
\qquad
r_M(c,t)=c
\]
for all \(m\in M\), \(c\in C\), and \(t\in I\).

Define
\[
H:Y\times I\to Y
\]
by
\[
H(y,t)=
\begin{cases}
	y,& y\in A,\\[4pt]
	r_M(y,t),& y\in M.
\end{cases}
\]
This is well-defined and continuous because on the overlap \(A\cap M=C\), the two formulas agree:
\[
r_M(c,t)=c.
\]

At time \(t=0\), we have
\[
H(y,0)=y
\]
for all \(y\in Y\), so \(H(-,0)=\mathrm{id}_Y\).

At time \(t=1\), every point of \(M\) has been moved into the common circle \(C\subset A\), while every point of \(A\) has remained fixed. Hence
\[
H(Y,1)=A.
\]
Thus \(H\) is a deformation retraction of \(Y\) onto the annulus \(A\).

\medskip

\textbf{Why \(Y\) deformation retracts onto the Möbius band.}
Exactly similarly, the annulus deformation retracts onto its core circle. Let
\[
r_A:A\times I\to A
\]
be a deformation retraction onto \(C\), fixing \(C\) pointwise.

Define
\[
K:Y\times I\to Y
\]
by
\[
K(y,t)=
\begin{cases}
	r_A(y,t),& y\in A,\\[4pt]
	y,& y\in M.
\end{cases}
\]
Again this is continuous because on the common circle \(C\),
\[
r_A(c,t)=c.
\]

At time \(t=0\), \(K\) is the identity on \(Y\). At time \(t=1\), all of \(A\) has collapsed into \(C\subset M\), while \(M\) has remained fixed. Therefore
\[
K(Y,1)=M.
\]
So \(K\) is a deformation retraction of \(Y\) onto the Möbius band \(M\).

\medskip

\textbf{Conclusion.}
The space obtained by gluing an annulus and a Möbius band along their core circles contains both of them as deformation retracts.

\subsection{Explicit coordinate models for the two bands}

It is often useful to write the annulus and the Möbius band explicitly as quotients.

The annulus may be written as
\[
A=[0,1]\times [-1,1]\big/\bigl((0,s)\sim (1,s)\bigr),
\]
and its core circle is
\[
C_A=[0,1]\times \{0\}\big/\bigl((0,0)\sim (1,0)\bigr).
\]
A deformation retraction of \(A\) onto \(C_A\) is given by
\[
r_A([u,s],t)=[u,(1-t)s].
\]

The Möbius band may be written as
\[
M=[0,1]\times [-1,1]\big/\bigl((0,s)\sim (1,-s)\bigr),
\]
and its core circle is
\[
C_M=[0,1]\times \{0\}\big/\bigl((0,0)\sim (1,0)\bigr).
\]
A deformation retraction of \(M\) onto \(C_M\) is given by
\[
r_M([u,s],t)=[u,(1-t)s].
\]

Since both core circles are naturally homeomorphic to \(S^1\), we may identify them and form the space \(Y\) above.	

\section{Theory Behind the Two Problems: Strong Deformation Retractions, Local Contractibility, and Gluing Along a Common Spine}

\subsection{Introduction}

The two problems considered above are closely related. At first sight they look different:

\begin{itemize}
	\item one is a local statement about neighbourhoods of a point in a contractible space,
	\item the other is a global construction of a space containing two different surfaces as deformation retracts.
\end{itemize}

However, both problems are governed by the same central idea:

\[
\text{a homotopy which fixes a distinguished subspace pointwise.}
\]

This is precisely the role of a \emph{strong deformation retraction}. Such homotopies are stronger than ordinary homotopies, and because they keep a chosen subset fixed at all times, they behave especially well in two situations:

\begin{enumerate}
	\item when we want to restrict a homotopy to smaller neighbourhoods,
	\item when we want to glue spaces together along a common subspace.
\end{enumerate}

Problem 5 uses the first principle. Problem 6 uses the second.

The goal of this section is to explain the theory behind both of them in a unified way.

\subsection{Basic definitions}

Let \(A\subset X\) be a subspace.

\subsubsection{Retraction}

A map
\[
r:X\to A
\]
is called a \emph{retraction} if
\[
r(a)=a
\qquad\text{for all }a\in A.
\]
Equivalently, \(r|_A=\mathrm{id}_A\).

In this case we say that \(A\) is a \emph{retract} of \(X\).

\subsubsection{Deformation retraction}

A \emph{deformation retraction} of \(X\) onto \(A\) is a homotopy
\[
H:X\times I\to X
\]
such that
\[
H(x,0)=x,
\qquad
H(x,1)\in A,
\qquad
H(a,1)=a
\quad\text{for all }a\in A.
\]

Thus \(H\) starts from the identity map on \(X\) and ends at a retraction of \(X\) onto \(A\).

If such a homotopy exists, we say that \(A\) is a \emph{deformation retract} of \(X\).

\subsubsection{Strong deformation retraction}

A \emph{strong deformation retraction} of \(X\) onto \(A\) is a deformation retraction
\[
H:X\times I\to X
\]
with the stronger condition
\[
H(a,t)=a
\qquad\text{for all }a\in A,\ t\in I.
\]

So the subspace \(A\) remains fixed pointwise during the entire homotopy, not just at the final time \(t=1\).

\medskip

\textbf{Why this matters.}
The difference between a deformation retract and a strong deformation retract is exactly what makes Problems 5 and 6 work so cleanly. If the distinguished subspace is fixed throughout the homotopy, then one can:

\begin{itemize}
	\item keep track of basepoints in pointed homotopy theory,
	\item glue homotopies on different pieces of a space,
	\item restrict homotopies to suitable neighbourhoods.
\end{itemize}

\subsection{Contractibility and pointed contractibility}

A space \(X\) is called \emph{contractible} if the identity map is homotopic to a constant map:
\[
\mathrm{id}_X \simeq c_{x_0}
\]
for some point \(x_0\in X\).

If there is a homotopy
\[
H:X\times I\to X
\]
such that
\[
H(x,0)=x,
\qquad
H(x,1)=x_0,
\qquad
H(x_0,t)=x_0
\]
for all \(x\in X\) and \(t\in I\), then \(X\) strongly deformation retracts onto the one-point subspace \(\{x_0\}\).

In that case, the contraction is compatible with the chosen basepoint. One may think of this as a pointed or based form of contractibility.

\medskip

\textbf{Remark.}
Every strong deformation retraction of \(X\) onto \(\{x_0\}\) makes \((X,x_0)\) trivial from the viewpoint of pointed homotopy theory. In particular, all based homotopy groups vanish:
\[
\pi_n(X,x_0)=0
\qquad\text{for all }n\geq 1.
\]

\subsection{Based maps and based homotopies}

A \emph{pointed space} is a pair \((X,x_0)\), where \(x_0\in X\) is the chosen basepoint.

A map
\[
f:(X,x_0)\to (Y,y_0)
\]
is called \emph{based} if
\[
f(x_0)=y_0.
\]

A \emph{based homotopy} between two based maps
\[
f,g:(X,x_0)\to (Y,y_0)
\]
is a homotopy
\[
F:X\times I\to Y
\]
such that
\[
F(x,0)=f(x),
\qquad
F(x,1)=g(x),
\qquad
F(x_0,t)=y_0
\quad\text{for all }t\in I.
\]

This last condition is essential: the basepoint must stay fixed throughout the homotopy.

\medskip

\textbf{Importance.}
Based homotopies are the correct notion in the study of:

\[
\pi_n(X,x_0),
\qquad
\Omega(X,x_0),
\qquad
\widetilde{H}_n(X),
\]
and many other pointed constructions.

Problem 5 is therefore naturally a \emph{based} statement, not just an unpointed one.

\subsection{The theory behind Problem 5: local pointed contractibility}

\subsubsection{Statement recalled}

Suppose \(X\) strongly deformation retracts to a point \(x_0\in X\). Then for every open neighbourhood \(U\) of \(x_0\), there exists a smaller open neighbourhood
\[
x_0\in V\subset U
\]
such that the inclusion
\[
(V,x_0)\hookrightarrow (U,x_0)
\]
is based homotopic to the constant map.

\subsubsection{Conceptual meaning}

This theorem says that a \emph{global} strong contraction of \(X\) to \(x_0\) implies a \emph{local} pointed contractibility property near \(x_0\).

The identity map of \(X\) is globally based homotopic to the constant map:
\[
\mathrm{id}_{(X,x_0)} \simeq c_{x_0}.
\]
The theorem says that this triviality can be seen already inside sufficiently small neighbourhoods: every neighbourhood \(U\) contains a smaller one \(V\) whose inclusion into \(U\) is based null-homotopic.

So the inclusion
\[
V\hookrightarrow U
\]
is, from the pointed homotopy viewpoint, homotopically trivial.

\subsubsection{Relation to local contractibility}

A space \(X\) is called \emph{locally contractible at \(x_0\)} if every neighbourhood \(U\ni x_0\) contains a smaller neighbourhood \(V\ni x_0\) which is contractible in \(U\). This means that the inclusion
\[
V\hookrightarrow U
\]
is homotopic to a constant map.

Problem 5 gives a stronger, pointed version of this statement. Namely, the homotopy can be chosen to preserve the basepoint throughout:
\[
(V,x_0)\hookrightarrow (U,x_0)
\]
is based null-homotopic.

Thus Problem 5 may be viewed as a theorem about \emph{pointed local contractibility}.

\subsubsection{Why the proof works}

The proof rests on two basic facts:

\begin{enumerate}
	\item continuity of the strong deformation retraction
	\[
	H:X\times I\to X,
	\]
	\item compactness of the interval
	\[
	I=[0,1].
	\]
\end{enumerate}

Since
\[
H(x_0,t)=x_0\in U
\qquad\text{for all }t\in I,
\]
continuity implies that for each fixed \(t\), points sufficiently near \(x_0\) also remain inside \(U\) for times near \(t\).

This gives an open neighbourhood of \((x_0,t)\) in \(X\times I\) which is sent into \(U\).

As \(t\) varies in \(I\), these neighbourhoods cover the compact set
\[
\{x_0\}\times I.
\]
A finite subcover then yields one smaller neighbourhood \(V\) of \(x_0\) which works simultaneously for all \(t\in I\). Hence
\[
H(V\times I)\subset U,
\]
and the restriction of \(H\) gives the desired based null-homotopy.

\subsubsection{General local principle}

The argument illustrates a general principle:

\medskip

\textbf{Local restriction principle.}
If a homotopy fixes a point \(x_0\) and stays inside a prescribed open set at that point, then sufficiently near \(x_0\) the homotopy also stays inside that open set.

\medskip

This is a standard continuity phenomenon, but when combined with compactness in the time variable, it becomes a powerful tool in homotopy theory.

\subsection{Examples illustrating Problem 5}

\subsubsection{Convex subsets of \(\mathbb{R}^n\)}

Let \(X\subset \mathbb{R}^n\) be convex, and let \(x_0\in X\). Then
\[
H(x,t)=(1-t)x+tx_0
\]
is a strong deformation retraction of \(X\) onto \(x_0\).

If \(U\) is any open neighbourhood of \(x_0\) in \(X\), then one may choose a smaller convex neighbourhood
\[
x_0\in V\subset U.
\]
For \(x\in V\), the straight line segment from \(x\) to \(x_0\) lies in \(U\), so the inclusion
\[
(V,x_0)\hookrightarrow (U,x_0)
\]
is based null-homotopic.

This is the simplest geometric model for the theorem.

\subsubsection{The interval}

Take
\[
X=[-1,1],
\qquad
x_0=0.
\]
The homotopy
\[
H(x,t)=(1-t)x
\]
strongly deformation retracts \(X\) onto \(0\).

Given any open neighbourhood \(U\) of \(0\), for instance
\[
U=(-\varepsilon,\varepsilon)\cap [-1,1],
\]
we may take
\[
V=\left(-\frac{\varepsilon}{2},\frac{\varepsilon}{2}\right)\cap [-1,1].
\]
Then the inclusion
\[
(V,0)\hookrightarrow (U,0)
\]
is based homotopic to the constant map via
\[
K(x,t)=(1-t)x.
\]

\subsubsection{A star-shaped domain}

Let \(X\subset \mathbb{R}^n\) be star-shaped with respect to \(x_0\). Then
\[
H(x,t)=(1-t)x+tx_0
\]
gives a strong deformation retraction of \(X\) onto \(x_0\).

Again, sufficiently small neighbourhoods of \(x_0\) contract inside any larger neighbourhood. So Problem 5 should be understood as a broad abstraction of this familiar geometric fact.

\subsection{The theory behind Problem 6: gluing along a common spine}

\subsubsection{Statement recalled}

Construct a space which contains both the annulus
\[
S^1\times I
\]
and the Möbius band as deformation retracts.

\subsubsection{The shared homotopy type}

Both the annulus and the Möbius band deformation retract onto a circle.

For the annulus, the retraction is onto its middle circle. For the Möbius band, the retraction is onto its core circle.

Therefore both spaces have the same homotopy type:
\[
S^1\times I \simeq S^1,
\qquad
\text{Möbius band} \simeq S^1.
\]

This common circle is the essential homotopy-theoretic data of both spaces. One often calls such a subspace a \emph{spine}.

\subsubsection{The spine idea}

A \emph{spine} of a space \(X\) is an embedded subspace \(C\subset X\) such that \(X\) deformation retracts onto \(C\). In many examples, \(C\) is much simpler than \(X\) but still captures its homotopy type.

For the annulus, the spine is a circle. For the Möbius band, the spine is also a circle.

Although these two surfaces are topologically different, they are homotopically the same because they collapse onto the same model space \(S^1\).

This suggests a natural construction: glue the two spaces along their common spine.

\subsection{A general gluing lemma}

The construction in Problem 6 is an instance of the following very useful fact.

\subsubsection{Lemma}

Let \(Y=A\cup B\), and suppose
\[
A\cap B=C.
\]
Assume that \(B\) strongly deformation retracts onto \(C\). Then \(Y\) deformation retracts onto \(A\).

\medskip

\textbf{Explanation.}
Since \(B\) strongly deformation retracts onto \(C\), there exists a homotopy
\[
r_B:B\times I\to B
\]
such that
\[
r_B(b,0)=b,
\qquad
r_B(b,1)\in C,
\qquad
r_B(c,t)=c
\quad\text{for all }c\in C.
\]

Now define
\[
H:Y\times I\to Y
\]
by
\[
H(y,t)=
\begin{cases}
	y,& y\in A,\\[4pt]
	r_B(y,t),& y\in B.
\end{cases}
\]

On the overlap \(C=A\cap B\), the two formulas agree because
\[
r_B(c,t)=c.
\]
Therefore \(H\) is well-defined and continuous.

At time \(t=0\), it is the identity on \(Y\). At time \(t=1\), all points of \(B\) have been pushed into
\[
C\subset A,
\]
while all points of \(A\) have remained fixed. Hence the image is exactly \(A\), so \(H\) is a deformation retraction of \(Y\) onto \(A\).

\medskip

\textbf{Symmetric version.}
If \(A\) strongly deformation retracts onto \(C\), then \(Y\) deformation retracts onto \(B\).

\medskip

This is exactly the principle used in Problem 6.

\subsection{Application to the annulus and the Möbius band}

Let \(A\) be an annulus and \(M\) a Möbius band. Let
\[
C_A\subset A,
\qquad
C_M\subset M
\]
be their core circles. Choose a homeomorphism
\[
\varphi:C_A\to C_M
\]
and identify these circles. Denote the resulting common image by
\[
C.
\]

Now form the quotient space
\[
Y=(A\sqcup M)/\sim,
\]
where \(c\sim \varphi(c)\) for all \(c\in C_A\).

Then
\[
Y=A\cup_C M.
\]

Since \(M\) strongly deformation retracts onto \(C\), the gluing lemma shows that \(Y\) deformation retracts onto \(A\).

Since \(A\) strongly deformation retracts onto \(C\), the symmetric version of the lemma shows that \(Y\) deformation retracts onto \(M\).

Thus \(Y\) contains both the annulus and the Möbius band as deformation retracts.

\subsection{Why this is interesting}

The annulus and the Möbius band are not homeomorphic:

\begin{itemize}
	\item the annulus has two boundary components,
	\item the Möbius band has one boundary component,
	\item the annulus is orientable,
	\item the Möbius band is non-orientable.
\end{itemize}

So the fact that both arise as deformation retracts of the same space does \emph{not} mean they are topologically the same. What it means is that they share the same homotopy type.

This is a valuable lesson:

\[
\text{homeomorphism type is much finer than homotopy type.}
\]

A deformation retract preserves homotopy type, but it does not preserve all topological structure.

\subsection{The annulus and the Möbius band as interval bundles over \(S^1\)}

There is also a geometric interpretation of Problem 6.

Both the annulus and the Möbius band are interval bundles over the circle \(S^1\):

\begin{itemize}
	\item the annulus is the trivial interval bundle
	\[
	S^1\times I,
	\]
	\item the Möbius band is the nontrivial interval bundle over \(S^1\).
\end{itemize}

In each case the zero section, or middle section, is a copy of \(S^1\), and the whole bundle contracts onto that section by collapsing fibres.

Thus the common circle is not accidental. It is the natural section over the base \(S^1\) in both cases.

From this viewpoint, Problem 6 says:

\begin{quote}
	Take the trivial and nontrivial interval bundle over \(S^1\), and glue them along their zero section.
\end{quote}

This gives a single space that can collapse either to the trivial bundle or to the twisted one.

\subsection{Explicit coordinate models}

It is useful to write the two surfaces concretely.

\subsubsection{The annulus}

The annulus may be represented as
\[
A=[0,1]\times [-1,1]\big/\bigl((0,s)\sim (1,s)\bigr).
\]

Its core circle is
\[
C_A=[0,1]\times \{0\}\big/\bigl((0,0)\sim (1,0)\bigr).
\]

A strong deformation retraction onto the core is
\[
r_A([u,s],t)=[u,(1-t)s].
\]

Indeed,
\[
r_A([u,s],0)=[u,s],
\qquad
r_A([u,s],1)=[u,0]\in C_A,
\]
and if \([u,0]\in C_A\), then
\[
r_A([u,0],t)=[u,0]
\]
for all \(t\).

\subsubsection{The Möbius band}

The Möbius band may be represented as
\[
M=[0,1]\times [-1,1]\big/\bigl((0,s)\sim (1,-s)\bigr).
\]

Its core circle is
\[
C_M=[0,1]\times \{0\}\big/\bigl((0,0)\sim (1,0)\bigr).
\]

Again the strong deformation retraction onto the core is
\[
r_M([u,s],t)=[u,(1-t)s].
\]

This is well-defined because when \(s\) is replaced by \(-s\) at the ends, multiplying by \((1-t)\) preserves the same identification.

Thus both surfaces collapse fibrewise onto a circle.

\subsection{Homotopy-theoretic interpretation}

Since both \(A\) and \(M\) deformation retract onto \(S^1\), they satisfy
\[
\pi_1(A)\cong \mathbb{Z},
\qquad
\pi_1(M)\cong \mathbb{Z},
\]
and for \(n\geq 2\),
\[
\pi_n(A)=0,
\qquad
\pi_n(M)=0.
\]

Likewise,
\[
H_0(A)\cong \mathbb{Z},
\qquad
H_1(A)\cong \mathbb{Z},
\qquad
H_n(A)=0\ \text{for }n\geq 2,
\]
and the same holds for the Möbius band:
\[
H_0(M)\cong \mathbb{Z},
\qquad
H_1(M)\cong \mathbb{Z},
\qquad
H_n(M)=0\ \text{for }n\geq 2.
\]

So from the viewpoint of algebraic topology, both spaces look like the circle. Their difference is geometric rather than homotopy-theoretic.

\subsection{A unifying perspective on both problems}

The two problems may now be compared directly.

\subsubsection{Problem 5}

There is a strong deformation retraction
\[
X\searrow \{x_0\}.
\]
Because the point \(x_0\) is fixed throughout the homotopy, continuity allows us to restrict the contraction to sufficiently small neighbourhoods of \(x_0\). Hence we obtain a local based null-homotopy.

\subsubsection{Problem 6}

There are strong deformation retractions
\[
A\searrow C,
\qquad
M\searrow C.
\]
Because the common circle \(C\) is fixed throughout these homotopies, they can be glued to form global deformation retractions of the larger glued space
\[
Y=A\cup_C M
\]
onto either \(A\) or \(M\).

\subsubsection{Common principle}

In both situations, the key mechanism is:

\[
\text{the homotopy fixes a chosen subspace pointwise.}
\]

That is exactly why strong deformation retracts are more useful than arbitrary homotopies when doing:

\begin{itemize}
	\item pointed topology,
	\item local arguments,
	\item gluing constructions.
\end{itemize}

\subsection{Two general lemmas to remember}

\subsubsection{Lemma A: local pointed triviality}

If \(X\) strongly deformation retracts onto a point \(x_0\), then for every open neighbourhood \(U\ni x_0\) there exists an open neighbourhood
\[
x_0\in V\subset U
\]
such that the inclusion
\[
(V,x_0)\hookrightarrow (U,x_0)
\]
is based null-homotopic.

\medskip

This is the abstract form of Problem 5.

\subsubsection{Lemma B: collapse one side of a union}

If
\[
Y=A\cup B,
\qquad
A\cap B=C,
\]
and \(B\) strongly deformation retracts onto \(C\), then \(Y\) deformation retracts onto \(A\).

\medskip

This is the abstract form of one half of Problem 6. The other half is obtained by symmetry.

\subsection{Further examples of the gluing principle}

The annulus and Möbius band are only one example. The same principle appears in many other situations.

\subsubsection{Two cylinders glued along a circle}

If two cylinders
\[
S^1\times I
\]
are glued along a common middle circle, then the resulting space deformation retracts onto either cylinder.

\subsubsection{A disk and a collar}

If a collar neighbourhood of the boundary of a disk is glued back along the boundary circle, the larger space still deformation retracts onto the disk, because the collar retracts onto the boundary circle.

\subsubsection{CW-complex intuition}

In CW-complex language, many spaces are built by attaching cells along a subspace. If an attached part collapses back onto the attaching region, then the whole space often deformation retracts onto the remaining part. Problem 6 is a geometric surface-level example of this general philosophy.

\subsection{Conclusion}

The theory behind the two problems may be summarized as follows.

\begin{enumerate}
	\item A strong deformation retraction is a homotopy that fixes a chosen subspace pointwise throughout time.
	\item This pointwise fixed condition makes the homotopy compatible with basepoints, restrictions, and gluing.
	\item In Problem 5, this yields a local theorem: every neighbourhood of the basepoint contains a smaller one whose inclusion is based null-homotopic.
	\item In Problem 6, this yields a global construction: two spaces with the same spine can be glued along that spine to form a larger space which collapses onto either one.
\end{enumerate}

Thus both problems are manifestations of one and the same principle:

\[
\text{strongly fixed subspaces allow controlled homotopy constructions.}
\]

This is one of the basic and most useful ideas in algebraic topology.


\section{Complements of Subspheres and Local Path Connectedness}

\subsection{Problem 7: The complement \(S^n\setminus S^m\)}

Let \(m<n\), and consider \(S^m\) as a subspace of \(S^n\) given by
\[
S^m
=
\left\{
(x_1,x_2,\dots,x_{m+1},0,\dots,0)
\;\middle|\;
\sum_{i=1}^{m+1} x_i^2 = 1
\right\}.
\]
We want to show that
\[
S^n \setminus S^m \simeq S^{\,n-m-1}.
\]

\paragraph{Intuition.}
Write
\[
\mathbb{R}^{n+1}=\mathbb{R}^{m+1}\times \mathbb{R}^{n-m}.
\]
Then a point of \(S^n\) can be written as \((u,v)\), where
\[
u\in \mathbb{R}^{m+1},
\qquad
v\in \mathbb{R}^{n-m},
\qquad
\|u\|^2+\|v\|^2=1.
\]
The embedded sphere \(S^m\) is exactly the set where \(v=0\). Thus its complement is the set of points with \(v\neq 0\). If \(v\neq 0\), then the direction of \(v\) determines a point of
\[
S^{n-m-1}\subset \mathbb{R}^{n-m},
\]
and one expects that the complement deformation retracts onto this sphere.

\paragraph{Examples.}
\begin{itemize}
	\item If \(m=0\), then \(S^0\subset S^n\) is just two antipodal points, and
	\[
	S^n\setminus S^0 \simeq S^{n-1}.
	\]
	
	\item If \(m=n-1\), then \(S^{n-1}\subset S^n\) is an equator, and its complement consists of two open hemispheres, so
	\[
	S^n\setminus S^{n-1}\simeq S^0.
	\]
	
	\item If \(n=3\) and \(m=1\), then
	\[
	S^3\setminus S^1 \simeq S^1.
	\]
\end{itemize}

\textbf{Solution.}
Let
\[
S^n
=
\left\{
(u,v)\in \mathbb{R}^{m+1}\times \mathbb{R}^{n-m}
\;\middle|\;
\|u\|^2+\|v\|^2=1
\right\}.
\]
Then
\[
S^m
=
\left\{
(u,0)\in S^n \;\middle|\; \|u\|=1
\right\},
\]
so
\[
S^n\setminus S^m
=
\left\{
(u,v)\in S^n \;\middle|\; v\neq 0
\right\}.
\]

Now define
\[
Y
=
\left\{
(0,v)\in S^n \;\middle|\; \|v\|=1
\right\}.
\]
This space is naturally homeomorphic to \(S^{n-m-1}\).

We define a homotopy
\[
H:(S^n\setminus S^m)\times [0,1]\to S^n\setminus S^m
\]
by
\[
H((u,v),t)
=
\frac{\bigl((1-t)u,\; v\bigr)}
{\sqrt{(1-t)^2\|u\|^2+\|v\|^2}}.
\]

We check that this is well defined.

Since \((u,v)\in S^n\setminus S^m\), we have \(v\neq 0\). Therefore
\[
(1-t)^2\|u\|^2+\|v\|^2 > 0,
\]
so the denominator never vanishes.

Also,
\[
\left\|
\frac{\bigl((1-t)u,\; v\bigr)}
{\sqrt{(1-t)^2\|u\|^2+\|v\|^2}}
\right\|^2
=
\frac{(1-t)^2\|u\|^2+\|v\|^2}
{(1-t)^2\|u\|^2+\|v\|^2}
=
1,
\]
so the image lies in \(S^n\). Since the second component remains a nonzero scalar multiple of \(v\), it is still nonzero, hence the image lies in \(S^n\setminus S^m\).

Now:
\[
H((u,v),0)
=
\frac{(u,v)}{\sqrt{\|u\|^2+\|v\|^2}}
=
(u,v),
\]
because \((u,v)\in S^n\).

At \(t=1\),
\[
H((u,v),1)
=
\frac{(0,v)}{\|v\|}
\in Y.
\]

Finally, if \((u,v)\in Y\), then \(u=0\) and \(\|v\|=1\), so
\[
H((0,v),t)
=
\frac{(0,v)}{\|v\|}
=
(0,v),
\]
which shows that every point of \(Y\) is fixed throughout the homotopy.

Therefore \(H\) is a strong deformation retraction of \(S^n\setminus S^m\) onto \(Y\). Since
\[
Y\cong S^{n-m-1},
\]
we conclude that
\[
S^n\setminus S^m \simeq S^{n-m-1}.
\]

\subsection{Problem 8: Connected and locally path connected implies path connected}

A space \(X\) is called \emph{locally path connected} if for every point \(x\in X\) and every neighbourhood \(U\ni x\), there exists a smaller neighbourhood \(V\) of \(x\) such that
\[
x\in V\subset U
\]
and \(V\) is path connected.

We want to prove:

\medskip

If \(X\) is connected and locally path connected, then \(X\) is path connected.

\paragraph{Intuition.}
Fix one point \(x_0\in X\), and consider the set of all points that can be joined to \(x_0\) by a path. Call this set \(P\). If the space is locally path connected, then small neighbourhoods can be used to extend paths, so \(P\) turns out to be both open and closed. Since \(X\) is connected and \(P\neq \varnothing\), we must have \(P=X\).

\paragraph{Why the hypothesis matters.}
Connected spaces need not be path connected. The classical counterexample is the topologist's sine curve. What fails there is local path connectedness. Thus local path connectedness is the extra condition that forces connectedness to imply path connectedness.

\textbf{Solution.}
Fix a point \(x_0\in X\), and define
\[
P
=
\left\{
x\in X
\;\middle|\;
\text{there exists a path from } x_0 \text{ to } x
\right\}.
\]
Clearly \(x_0\in P\), so \(P\neq \varnothing\).

We first show that \(P\) is open.

Let \(y\in P\). Since \(X\) is locally path connected, there exists a path connected neighbourhood \(V\) of \(y\). Because \(V\) is a neighbourhood of \(y\), there exists an open set \(O\) such that
\[
y\in O\subset V.
\]
Now take any \(z\in O\). Since both \(y\) and \(z\) lie in the path connected set \(V\), there exists a path in \(V\) from \(y\) to \(z\). Since \(y\in P\), there is already a path from \(x_0\) to \(y\). Concatenating these two paths gives a path from \(x_0\) to \(z\). Hence \(z\in P\), so
\[
O\subset P.
\]
Therefore \(P\) is open.

Next we show that \(P\) is closed.

Let \(y\in \overline{P}\). Again, by local path connectedness, there exists a path connected neighbourhood \(V\) of \(y\). Since \(y\in \overline{P}\), every neighbourhood of \(y\) meets \(P\), so in particular
\[
V\cap P\neq \varnothing.
\]
Choose a point
\[
z\in V\cap P.
\]
Because \(z\in P\), there exists a path from \(x_0\) to \(z\). Since \(V\) is path connected and \(y,z\in V\), there exists a path from \(z\) to \(y\). Concatenating these paths yields a path from \(x_0\) to \(y\). Thus \(y\in P\).

So
\[
\overline{P}\subset P,
\]
which means that \(P\) is closed.

We have shown that \(P\) is nonempty, open, and closed in \(X\). Since \(X\) is connected, the only such subset is the whole space. Therefore
\[
P=X.
\]
Hence every point of \(X\) can be joined to \(x_0\) by a path, and so \(X\) is path connected.

\paragraph{Conclusion.}
A connected locally path connected space has exactly one path component, so it is path connected.

\section{Path Connectedness, Local Path Connectedness, and Counterexamples}

\subsection{Introduction}

The notions of connectedness and path connectedness are among the most basic and most important global properties in topology. Although they are closely related, they are not the same. A space may be connected without admitting paths between all pairs of points. The extra local hypothesis of local path connectedness repairs this defect and yields one of the standard theorems of general topology:
\[
\text{connected} + \text{locally path connected} \implies \text{path connected}.
\]

This section develops the theory carefully. We begin with the definitions, then explain the relations among these notions, then discuss connected components and path components, and finally present the standard counterexamples.

\subsection{Basic definitions}

\subsubsection{Connectedness}

A topological space \(X\) is called \emph{connected} if it cannot be written as a union
\[
X=U\cup V
\]
where \(U\) and \(V\) are disjoint nonempty open subsets of \(X\).

Equivalently, \(X\) is connected if the only subsets of \(X\) that are both open and closed are
\[
\varnothing
\qquad\text{and}\qquad
X.
\]

Another equivalent formulation is that every continuous map
\[
f:X\to \{0,1\}
\]
to a discrete two-point space is constant.

\paragraph{Intuition.}
Connectedness means that the space is ``in one piece'' from the topological point of view. However, connectedness alone does not guarantee that one can continuously move inside the space from any point to any other point.

\subsubsection{Paths and path connectedness}

A \emph{path} in a space \(X\) is a continuous map
\[
\gamma:[0,1]\to X.
\]
If
\[
\gamma(0)=x
\qquad\text{and}\qquad
\gamma(1)=y,
\]
then \(\gamma\) is called a path from \(x\) to \(y\).

A space \(X\) is called \emph{path connected} if for every pair of points \(x,y\in X\), there exists a path
\[
\gamma:[0,1]\to X
\]
with
\[
\gamma(0)=x,
\qquad
\gamma(1)=y.
\]

\paragraph{Intuition.}
Path connectedness is stronger than connectedness. It says not merely that the space is not broken into disjoint open pieces, but that any two points can actually be joined by a continuous curve lying entirely inside the space.

\subsubsection{Local connectedness}

A space \(X\) is called \emph{locally connected} if for every point \(x\in X\) and every neighbourhood \(U\ni x\), there exists a connected neighbourhood \(V\) such that
\[
x\in V\subset U.
\]

Equivalently, every point has a neighbourhood basis consisting of connected open sets.

\subsubsection{Local path connectedness}

A space \(X\) is called \emph{locally path connected} if for every point \(x\in X\) and every neighbourhood \(U\ni x\), there exists a path connected neighbourhood \(V\) such that
\[
x\in V\subset U.
\]

Equivalently, every point has a neighbourhood basis consisting of path connected open sets.

\paragraph{Immediate observation.}
Since every path connected space is connected, every path connected neighbourhood is connected. Therefore
\[
\text{locally path connected} \implies \text{locally connected}.
\]
The converse is false in general.

\subsection{The basic implication: path connected implies connected}

\subsubsection{Theorem}

If \(X\) is path connected, then \(X\) is connected.

\paragraph{\textbf{Solution.}}
Assume that \(X\) is path connected but not connected. Then there exist disjoint nonempty open subsets \(U,V\subset X\) such that
\[
X=U\cup V.
\]
Choose points
\[
x\in U,
\qquad
y\in V.
\]
Since \(X\) is path connected, there exists a path
\[
\gamma:[0,1]\to X
\]
with
\[
\gamma(0)=x,
\qquad
\gamma(1)=y.
\]

Now consider the sets
\[
\gamma^{-1}(U)
\qquad\text{and}\qquad
\gamma^{-1}(V).
\]
These are open in \([0,1]\) because \(\gamma\) is continuous and \(U,V\) are open. They are disjoint, and their union is all of \([0,1]\), because
\[
X=U\cup V.
\]
Also,
\[
0\in \gamma^{-1}(U),
\qquad
1\in \gamma^{-1}(V),
\]
so both sets are nonempty.

Thus \([0,1]\) is written as a union of two disjoint nonempty open subsets, contradicting the fact that \([0,1]\) is connected.

Therefore \(X\) must be connected.

\paragraph{Conclusion.}
We have proved
\[
\text{path connected} \implies \text{connected}.
\]

\subsection{Why the converse fails}

The converse
\[
\text{connected} \implies \text{path connected}
\]
is false in general.

The classical counterexample is the \emph{topologist's sine curve}, which we discuss later in detail.

Thus the difference between connectedness and path connectedness is real and important.

\subsection{Connected components and path components}

\subsubsection{Connected components}

A \emph{connected component} of \(X\) is a maximal connected subset of \(X\), that is, a connected subset that is not properly contained in any larger connected subset.

Connected components always exist, and every point of \(X\) belongs to exactly one connected component. Hence \(X\) is the disjoint union of its connected components.

\subsubsection{Path components}

A \emph{path component} of \(X\) is a maximal path connected subset of \(X\).

Equivalently, for a point \(x\in X\), the path component of \(x\) is
\[
P(x)=\{y\in X : \text{there exists a path from } x \text{ to } y\}.
\]

Again, every point belongs to exactly one path component, and \(X\) is the disjoint union of its path components.

\subsubsection{Relation between the two}

Since path connected implies connected, every path component is connected. Therefore every path component is contained in some connected component.

In symbols:
\[
\text{path component of }x
\subset
\text{connected component of }x.
\]

In general, this inclusion may be strict. A connected component can contain many different path components.

\paragraph{Important idea.}
The failure of connectedness to imply path connectedness means precisely that a connected component need not be path connected.

\subsection{Local path connectedness and openness of path components}

One of the most important facts is that in a locally path connected space, path components are open.

\subsubsection{Theorem}

Let \(X\) be locally path connected. Then every path component of \(X\) is open.

\paragraph{\textbf{Solution.}}
Let \(P\) be a path component of \(X\), and let \(x\in P\). Since \(X\) is locally path connected, for every neighbourhood \(U\) of \(x\), there exists a path connected neighbourhood \(V\) such that
\[
x\in V\subset U.
\]

In particular, choose any open neighbourhood \(U\) of \(x\), and let \(V\) be an open path connected neighbourhood with
\[
x\in V\subset U.
\]
Because \(V\) is path connected and contains \(x\), every point of \(V\) can be joined to \(x\) by a path lying in \(V\subset X\). Hence every point of \(V\) belongs to the same path component as \(x\), namely \(P\). Therefore
\[
V\subset P.
\]

So \(x\) has an open neighbourhood contained in \(P\). Since \(x\in P\) was arbitrary, \(P\) is open.

\paragraph{Conclusion.}
In locally path connected spaces, path components are open subsets.

\subsubsection{Corollary}

If \(X\) is locally path connected, then connected components and path components coincide.

\paragraph{\textbf{Solution.}}
Let \(C\) be a connected component of \(X\). Since path components partition \(X\), the set \(C\) is the union of all path components contained in it. But in a locally path connected space, path components are open. Hence \(C\) is written as a union of disjoint open subsets of the subspace \(C\), namely its path components.

Since \(C\) is connected, it cannot be the union of more than one disjoint nonempty open subset. Therefore \(C\) must consist of exactly one path component.

Thus every connected component is path connected, and since every path component is always contained in a connected component, the two notions coincide.

\paragraph{Conclusion.}
For locally path connected spaces,
\[
\text{connected component}=\text{path component}.
\]

\subsection{Connected plus locally path connected implies path connected}

We now prove the standard theorem used repeatedly in topology.

\subsubsection{Theorem}

If \(X\) is connected and locally path connected, then \(X\) is path connected.

\paragraph{\textbf{Solution.}}
Fix a point \(x_0\in X\), and define
\[
P
=
\{x\in X : \text{there exists a path from }x_0\text{ to }x\}.
\]
This is the path component of \(x_0\).

Clearly,
\[
x_0\in P,
\]
so \(P\neq \varnothing\).

We claim that \(P\) is open. Indeed, since \(X\) is locally path connected, every path component is open, by the theorem above. Hence \(P\) is open.

We also claim that \(P\) is closed. There are two ways to see this.

\paragraph{First proof of closedness.}
Because path components coincide with connected components in a locally path connected space, \(P\) is a connected component. Connected components are always closed, so \(P\) is closed.

\paragraph{Direct proof of closedness.}
Let
\[
y\in \overline{P}.
\]
Since \(X\) is locally path connected, there exists an open path connected neighbourhood \(V\) of \(y\). Because
\[
y\in \overline{P},
\]
the neighbourhood \(V\) meets \(P\), so choose a point
\[
z\in V\cap P.
\]
Since \(z\in P\), there is a path from \(x_0\) to \(z\). Since \(V\) is path connected and \(y,z\in V\), there is a path from \(z\) to \(y\). Concatenating these two paths gives a path from \(x_0\) to \(y\). Hence
\[
y\in P.
\]
Thus
\[
\overline{P}\subset P,
\]
so \(P\) is closed.

Now \(P\) is nonempty, open, and closed in \(X\). Since \(X\) is connected, the only nonempty subset with these properties is all of \(X\). Therefore
\[
P=X.
\]

Thus every point of \(X\) can be joined to \(x_0\) by a path, so \(X\) is path connected.

\paragraph{Conclusion.}
We have proved
\[
\text{connected}+\text{locally path connected}\implies \text{path connected}.
\]

\subsection{Examples of spaces that are path connected and locally path connected}

\subsubsection{\(\mathbb{R}^n\)}

The Euclidean space \(\mathbb{R}^n\) is path connected and locally path connected.

Indeed, any two points \(a,b\in \mathbb{R}^n\) are joined by the line segment
\[
\gamma(t)=(1-t)a+tb,
\qquad
0\leq t\leq 1.
\]
Also, every open ball in \(\mathbb{R}^n\) is convex, hence path connected. So \(\mathbb{R}^n\) is locally path connected.

\subsubsection{Open subsets of \(\mathbb{R}^n\)}

Every open subset \(U\subset \mathbb{R}^n\) is locally path connected.

Indeed, if \(x\in U\), then because \(U\) is open there exists an \(\varepsilon>0\) such that the open ball
\[
B(x,\varepsilon)\subset U.
\]
This ball is convex, hence path connected. Therefore every point has a basis of path connected neighbourhoods.

An open subset of \(\mathbb{R}^n\) need not be connected. However, each of its connected components is automatically path connected by the theorem above.

\subsubsection{Spheres}

For \(n\geq 1\), the sphere
\[
S^n
\]
is path connected and locally path connected.

Indeed, locally it looks like Euclidean space, and globally any two points on \(S^n\) can be joined by an arc.

By contrast,
\[
S^0=\{-1,1\}
\]
is not connected and therefore not path connected.

\subsubsection{Manifolds}

Every topological manifold is locally path connected, because each point has a neighbourhood homeomorphic to an open subset of some \(\mathbb{R}^n\), and open subsets of \(\mathbb{R}^n\) are locally path connected.

Hence for manifolds,
\[
\text{connected} \iff \text{path connected}.
\]

\subsubsection{CW-complexes and simplicial complexes}

CW-complexes and simplicial complexes are also locally path connected. Therefore in these categories connectedness and path connectedness agree.

\subsection{The topologist's sine curve}

We now discuss the standard counterexample showing that connectedness does not imply path connectedness.

\subsubsection{Definition}

Consider the subset
\[
S
=
\left\{
\left(x,\sin\frac{1}{x}\right):0<x\leq 1
\right\}
\subset \mathbb{R}^2.
\]
Its closure in \(\mathbb{R}^2\) is
\[
T
=
\overline{S}
=
\left\{
\left(x,\sin\frac{1}{x}\right):0<x\leq 1
\right\}
\cup
\bigl(\{0\}\times[-1,1]\bigr).
\]
The space \(T\) is called the \emph{topologist's sine curve}.

\subsubsection{Why \(T\) is connected}

The interval
\[
(0,1]
\]
is connected, and the map
\[
f:(0,1]\to \mathbb{R}^2,
\qquad
f(x)=\left(x,\sin\frac{1}{x}\right)
\]
is continuous. Therefore its image \(S\) is connected.

Now \(T=\overline{S}\) is the closure of a connected set. The closure of a connected set is connected. Hence \(T\) is connected.

\subsubsection{Why \(T\) is not path connected}

Although \(T\) is connected, it is not path connected. Intuitively, the oscillating graph
\[
y=\sin\frac{1}{x}
\]
wiggles infinitely often as \(x\to 0^+\), accumulating on the whole vertical segment
\[
\{0\}\times[-1,1].
\]
There is no continuous path in \(T\) joining a point on the vertical segment to a point on the oscillating part.

For example, there is no path in \(T\) joining
\[
(0,0)
\]
to
\[
\left(\frac{1}{\pi},0\right).
\]

A full proof usually proceeds by assuming such a path exists and then using continuity to show that the first coordinate of the path would have to cross infinitely many oscillations in an impossible way. The crucial fact is that the path would force the image of a connected interval to meet infinitely many alternating arcs, violating continuity properties near the first time the path leaves the vertical segment.

Thus \(T\) is connected but not path connected.

\paragraph{Conclusion.}
The topologist's sine curve proves that
\[
\text{connected} \centernot\implies \text{path connected}.
\]

\subsubsection{Why \(T\) is not locally path connected}

At points of the vertical segment
\[
\{0\}\times[-1,1],
\]
every sufficiently small neighbourhood intersects infinitely many oscillations of the graph. No matter how small the neighbourhood is, one cannot find a smaller open neighbourhood that is path connected in the subspace topology.

Thus \(T\) is not locally path connected.

This is exactly why the theorem
\[
\text{connected}+\text{locally path connected}\implies \text{path connected}
\]
does not apply to \(T\).

\subsection{Comb spaces}

Comb spaces are useful for intuition because they exhibit subtle local behavior.

\subsubsection{The comb space}

Define
\[
C
=
\bigl([0,1]\times\{0\}\bigr)
\cup
\bigl(\{0\}\times[0,1]\bigr)
\cup
\bigcup_{n=1}^{\infty}
\left(
\left\{\frac{1}{n}\right\}\times[0,1]
\right)
\subset \mathbb{R}^2.
\]

This space consists of:
\begin{itemize}
	\item the horizontal base segment \([0,1]\times\{0\}\),
	\item the left vertical spine \(\{0\}\times[0,1]\),
	\item and infinitely many vertical ``teeth'' at \(x=1/n\).
\end{itemize}

\subsubsection{Properties of the comb space}

The comb space \(C\) is connected.

In fact, \(C\) is also path connected. For example, given any point in the space, one can move vertically down to the base segment and then move horizontally along the base segment, and then move vertically if needed.

However, the local behavior near the spine \(\{0\}\times[0,1]\) is subtle, and comb-type examples are often modified to produce spaces that fail local connectedness or local path connectedness.

\subsubsection{Deleted comb space}

A useful variation is obtained by deleting part of the spine. For instance,
\[
C'
=
C\setminus \bigl(\{0\}\times(0,1]\bigr).
\]
Then \(C'\) has very different connectivity properties and is often used as a source of counterexamples in point-set topology.

These examples show that small-scale topology can behave in complicated ways even for subsets of the plane.

\subsection{Local connectedness versus local path connectedness}

We already know that
\[
\text{locally path connected} \implies \text{locally connected}.
\]
The converse is false.

That is, there exist spaces that have arbitrarily small connected neighbourhoods but do not have arbitrarily small path connected neighbourhoods.

Such examples are more delicate than the topologist's sine curve and often arise from specially constructed planar continua. For many elementary courses, it is enough to remember the logical distinction and the fact that local path connectedness is strictly stronger.

\subsection{Simple examples and nonexamples}

\subsubsection{Discrete spaces}

If \(X\) is a discrete space, then every singleton \(\{x\}\) is open and path connected. Hence every discrete space is locally path connected.

However, a discrete space is connected if and only if it has only one point. Thus a discrete two-point space is locally path connected but not connected and not path connected.

\subsubsection{Intervals}

Every interval in \(\mathbb{R}\) is path connected. A path from \(a\) to \(b\) is given by
\[
\gamma(t)=(1-t)a+tb.
\]
Hence every interval is connected.

This is an important source of examples because many connectedness arguments are reduced to statements about intervals.

\subsubsection{Union of disjoint intervals}

The space
\[
[0,1]\cup [2,3]\subset \mathbb{R}
\]
is not connected, hence not path connected. It is locally path connected, because open subsets of \(\mathbb{R}\) are locally path connected.

This shows that local path connectedness alone does not imply connectedness or path connectedness.

\subsubsection{Circle and torus}

The circle \(S^1\) and the torus \(T^2\) are path connected and locally path connected. They are standard examples where these notions behave exactly as expected.

\subsection{Logical summary of implications}

We summarize the main logical relations.

\subsubsection{Always true}

For every topological space,
\[
\text{path connected} \implies \text{connected},
\]
and
\[
\text{locally path connected} \implies \text{locally connected}.
\]

\subsubsection{True with an extra hypothesis}

If a space is locally path connected, then
\[
\text{connected} \implies \text{path connected}.
\]

Equivalently, in a locally path connected space,
\[
\text{connected} \iff \text{path connected}.
\]

\subsubsection{False in general}

The following converse implications fail in general:
\[
\text{connected} \centernot\implies \text{path connected},
\]
\[
\text{locally connected} \centernot\implies \text{locally path connected}.
\]

\subsection{A useful theorem about closures of connected sets}

Since this fact is used in the topologist's sine curve example, we record it separately.

\subsubsection{Theorem}

If \(A\subset X\) is connected, then its closure \(\overline{A}\) is connected.

\paragraph{\textbf{Solution.}}
Assume \(\overline{A}\) is not connected. Then there exist disjoint nonempty open subsets \(U,V\) of the subspace \(\overline{A}\) such that
\[
\overline{A}=U\cup V.
\]
Since \(A\subset \overline{A}\), the set \(A\) is partitioned as
\[
A=(A\cap U)\cup (A\cap V).
\]
The sets \(A\cap U\) and \(A\cap V\) are open in the subspace topology on \(A\), disjoint, and cover \(A\).

Since \(A\) is connected, one of them must be empty. Suppose
\[
A\cap U=\varnothing.
\]
Then
\[
A\subset V.
\]
Because \(V\) is closed in \(\overline{A}\) relative to the decomposition into clopen sets, it follows that \(\overline{A}\subset V\), so
\[
U=\varnothing,
\]
a contradiction.

Thus \(\overline{A}\) is connected.

\subsection{A useful theorem about continuous images}

We also record the basic fact used repeatedly above.

\subsubsection{Theorem}

The continuous image of a connected space is connected.

\paragraph{\textbf{Solution.}}
Let
\[
f:X\to Y
\]
be continuous, and suppose \(X\) is connected. Assume for contradiction that \(f(X)\) is disconnected. Then
\[
f(X)=U\cup V
\]
where \(U\) and \(V\) are disjoint nonempty open subsets of the subspace \(f(X)\).

Then
\[
f^{-1}(U)
\qquad\text{and}\qquad
f^{-1}(V)
\]
are disjoint nonempty open subsets of \(X\), and
\[
X=f^{-1}(U)\cup f^{-1}(V),
\]
contradicting connectedness of \(X\).

Hence \(f(X)\) is connected.

\subsection{What one should remember}

The most important points are the following.

\begin{enumerate}
	\item A path connected space is always connected.
	\item A connected space need not be path connected.
	\item The topologist's sine curve is the standard counterexample.
	\item Local path connectedness is the local condition that makes connectedness imply path connectedness.
	\item In a locally path connected space, path components are open.
	\item In a locally path connected space, connected components and path components coincide.
\end{enumerate}

\subsection{Compact final summary}

The theory may be compressed into the following diagram:
\[
\text{path connected}
\implies
\text{connected},
\]
\[
\text{locally path connected}
\implies
\text{locally connected}.
\]

Neither converse is true in general.

However, under local path connectedness one obtains the crucial theorem:
\[
\text{connected}+\text{locally path connected}
\implies
\text{path connected}.
\]

Thus for spaces such as manifolds, CW-complexes, open subsets of Euclidean space, and many geometric objects of interest, connectedness and path connectedness are equivalent notions. The standard obstruction comes from spaces with complicated small-scale oscillatory behavior, of which the topologist's sine curve is the classical example.


\subsection{Implication diagram, examples, and counterexamples}

It is tempting to guess a chain
\[
\text{connected} \implies \text{path connected} \implies \text{locally path connected},
\]
but this is false. The correct picture is more subtle.

\subsubsection{Clean implication diagram}

The always valid implication is
\[
\text{path connected} \implies \text{connected}.
\]

A second important theorem is
\[
\text{connected}+\text{locally path connected} \implies \text{path connected}.
\]

Thus the correct logical picture is:

\[
\begin{array}{c}
	\text{path connected} \Longrightarrow \text{connected},\\[8pt]
	\text{connected}+\text{locally path connected} \Longrightarrow \text{path connected}.
\end{array}
\]

The following implications are false in general:
\[
\text{connected} \not\Rightarrow \text{path connected},
\]
\[
\text{path connected} \not\Rightarrow \text{locally path connected},
\]
\[
\text{connected} \not\Rightarrow \text{locally path connected},
\]
\[
\text{locally path connected} \not\Rightarrow \text{connected},
\]
and therefore also
\[
\text{locally path connected} \not\Rightarrow \text{path connected}.
\]

\subsubsection{A compact summary table}

\[
\begin{array}{|c|c|c|}
	\hline
	\text{Statement} & \text{True or false?} & \text{Standard example or counterexample} \\
	\hline
	\text{path connected} \implies \text{connected} & \text{true} & \mathbb{R}^n,\ S^n,\ [0,1] \\
	\hline
	\text{connected} \implies \text{path connected} & \text{false} & \text{topologist's sine curve} \\
	\hline
	\text{path connected} \implies \text{locally path connected} & \text{false} & \text{comb space} \\
	\hline
	\text{locally path connected} \implies \text{connected} & \text{false} & [0,1]\cup[2,3] \\
	\hline
	\text{locally path connected} \implies \text{path connected} & \text{false} & [0,1]\cup[2,3] \\
	\hline
	\text{connected} \implies \text{locally path connected} & \text{false} & \text{topologist's sine curve} \\
	\hline
\end{array}
\]

\subsubsection{Example 1: a space with all three properties}

Let
\[
X=\mathbb{R}^n.
\]

Then \(X\) is path connected, because for any \(a,b\in \mathbb{R}^n\), the map
\[
\gamma:[0,1]\to \mathbb{R}^n,
\qquad
\gamma(t)=(1-t)a+tb,
\]
is continuous and satisfies
\[
\gamma(0)=a,
\qquad
\gamma(1)=b.
\]
Thus \(\mathbb{R}^n\) is path connected.

Since path connected spaces are connected, \(\mathbb{R}^n\) is connected.

Also, \(\mathbb{R}^n\) is locally path connected, because every point \(x\in \mathbb{R}^n\) has arbitrarily small open balls
\[
B(x,\varepsilon)
\]
as neighbourhoods, and each open ball is convex, hence path connected.

Therefore \(\mathbb{R}^n\) is:
\[
\text{connected},\qquad \text{path connected},\qquad \text{locally path connected}.
\]

Exactly the same conclusion holds for:
\[
S^n \text{ for } n\geq 1,\qquad [0,1],\qquad \text{open balls},\qquad \text{connected manifolds}.
\]

\subsubsection{Example 2: locally path connected but not connected, hence not path connected}

Let
\[
X=[0,1]\cup[2,3]\subset \mathbb{R}.
\]

We first show that \(X\) is locally path connected.

Indeed, if \(x\in X\), then there exists a sufficiently small open interval
\[
(x-\varepsilon,x+\varepsilon)
\]
in \(\mathbb{R}\) such that
\[
(x-\varepsilon,x+\varepsilon)\cap X
\]
is either an interval in \([0,1]\) or an interval in \([2,3]\). Any interval is path connected, so each point has arbitrarily small path connected neighbourhoods in the subspace topology. Hence \(X\) is locally path connected.

However, \(X\) is not connected, because
\[
[0,1]
\qquad\text{and}\qquad
[2,3]
\]
are disjoint nonempty open-and-closed subsets in the subspace topology on \(X\), and
\[
X=[0,1]\cup[2,3].
\]

Since path connected spaces are connected, \(X\) cannot be path connected.

Thus \(X\) is an example of a space that is
\[
\text{locally path connected but not connected,}
\]
and therefore also
\[
\text{locally path connected but not path connected.}
\]

\subsubsection{Example 3: connected but not path connected, and not locally path connected}

Define
\[
S=
\left\{
\left(x,\sin\frac{1}{x}\right):0<x\leq 1
\right\}\subset \mathbb{R}^2
\]
and let
\[
T=\overline{S}
=
\left\{
\left(x,\sin\frac{1}{x}\right):0<x\leq 1
\right\}
\cup
\bigl(\{0\}\times[-1,1]\bigr).
\]
This space \(T\) is the topologist's sine curve.

\paragraph{Claim 1.}
\(T\) is connected.

\paragraph{Proof.}
The interval
\[
(0,1]
\]
is connected, and the map
\[
f:(0,1]\to \mathbb{R}^2,
\qquad
f(x)=\left(x,\sin\frac{1}{x}\right),
\]
is continuous. Hence its image \(S\) is connected. Since the closure of a connected set is connected, \(\overline{S}=T\) is connected.

\paragraph{Claim 2.}
\(T\) is not path connected.

\paragraph{Explanation.}
The oscillating part
\[
\left\{
\left(x,\sin\frac{1}{x}\right):0<x\leq 1
\right\}
\]
accumulates on the whole vertical segment
\[
\{0\}\times[-1,1].
\]
There is no path in \(T\) joining a point on the vertical segment, such as \((0,0)\), to a point on the oscillating graph, such as
\[
\left(\frac{1}{\pi},0\right).
\]
This is the classical example of a connected but not path connected space.

\paragraph{Claim 3.}
\(T\) is not locally path connected.

\paragraph{Explanation.}
At any point of the vertical segment
\[
\{0\}\times[-1,1],
\]
every neighbourhood in the subspace topology meets infinitely many oscillations of the graph. One cannot find arbitrarily small path connected neighbourhoods there. Hence \(T\) is not locally path connected.

Therefore \(T\) is an example of a space that is
\[
\text{connected but not path connected},
\]
and also
\[
\text{connected but not locally path connected}.
\]

\subsubsection{Example 4: path connected but not locally path connected}

Define the comb space
\[
C=
\bigl([0,1]\times\{0\}\bigr)
\cup
\bigl(\{0\}\times[0,1]\bigr)
\cup
\bigcup_{n=1}^{\infty}
\left(
\left\{\frac{1}{n}\right\}\times[0,1]
\right)
\subset \mathbb{R}^2.
\]

This consists of:
\begin{itemize}
	\item the horizontal base segment \([0,1]\times\{0\}\),
	\item the vertical spine \(\{0\}\times[0,1]\),
	\item and the infinitely many vertical teeth at \(x=1/n\).
\end{itemize}

\paragraph{Claim 1.}
\(C\) is path connected.

\paragraph{Proof.}
Let
\[
p=(x_1,y_1)\in C,
\qquad
q=(x_2,y_2)\in C.
\]
We construct a path from \(p\) to \(q\) in three steps.

First, move vertically from \(p\) to the base:
\[
\gamma_1(t)=(x_1,(1-t)y_1),
\qquad
0\leq t\leq 1.
\]
Because \(p\in C\), the point \(x_1\) is either \(0\) or \(1/n\) for some \(n\), so this vertical segment lies in \(C\).

Second, move horizontally along the base from \((x_1,0)\) to \((x_2,0)\):
\[
\gamma_2(t)=((1-t)x_1+tx_2,0),
\qquad
0\leq t\leq 1.
\]
This lies in \([0,1]\times\{0\}\subset C\).

Third, move vertically from \((x_2,0)\) to \(q\):
\[
\gamma_3(t)=(x_2,ty_2),
\qquad
0\leq t\leq 1.
\]
Again this lies in \(C\).

By concatenating \(\gamma_1\), \(\gamma_2\), and \(\gamma_3\), we obtain a path from \(p\) to \(q\). Thus \(C\) is path connected.

\paragraph{Claim 2.}
\(C\) is not locally path connected.

\paragraph{Explanation.}
Consider a point
\[
(0,y)\in \{0\}\times(0,1].
\]
Every sufficiently small neighbourhood of \((0,y)\) in the subspace topology intersects not only a short segment of the spine, but also infinitely many nearby teeth
\[
\left\{\frac{1}{n}\right\}\times[0,1].
\]
If the neighbourhood is chosen small enough so that it does not reach the base line
\[
[0,1]\times\{0\},
\]
then those tooth fragments are disconnected from the spine within that neighbourhood. Hence no sufficiently small neighbourhood of \((0,y)\) contains a smaller open path connected neighbourhood. Therefore \(C\) is not locally path connected at such points.

So \(C\) is a standard example of a space that is
\[
\text{path connected but not locally path connected.}
\]

Since path connected spaces are connected, \(C\) is also connected.

\subsubsection{Formal conclusion}

The examples above establish the exact logical situation:

\begin{itemize}
	\item \(\mathbb{R}^n\) shows that a space may have all three properties.
	\item \([0,1]\cup[2,3]\) shows that locally path connected does not imply connected, and does not imply path connected.
	\item The topologist's sine curve shows that connected does not imply path connected, and does not imply locally path connected.
	\item The comb space shows that path connected does not imply locally path connected.
\end{itemize}

Thus there is no linear chain
\[
\text{connected} \implies \text{path connected} \implies \text{locally path connected}.
\]

Instead, one must remember the correct facts:
\[
\text{path connected} \implies \text{connected},
\]
and
\[
\text{connected}+\text{locally path connected} \implies \text{path connected}.
\]

\subsubsection{A final diagram}

A clean way to remember the situation is the following:

\[
\begin{array}{c}
	\text{path connected} \Longrightarrow \text{connected},\\[10pt]
	\text{connected}+\text{locally path connected} \Longrightarrow \text{path connected},\\[10pt]
	\text{connected} \not\Rightarrow \text{path connected},\\[6pt]
	\text{path connected} \not\Rightarrow \text{locally path connected},\\[6pt]
	\text{connected} \not\Rightarrow \text{locally path connected},\\[6pt]
	\text{locally path connected} \not\Rightarrow \text{connected}.
\end{array}
\]

\subsection{Geometric meaning of \(S^n \setminus S^m \simeq S^{\,n-m-1}\): examples and intuition}

The statement
\[
S^n \setminus S^m \simeq S^{\,n-m-1}
\]
does \emph{not} say that the complement is literally equal to a sphere, and it does \emph{not} say that it is a quotient space. Rather, it says that the complement of the embedded sphere \(S^m\) inside \(S^n\) has the same \emph{homotopy type} as a sphere of dimension \(n-m-1\).

In other words, after continuously deforming the complement without tearing or gluing, one can shrink it onto a copy of
\[
S^{\,n-m-1}.
\]

\subsubsection{What the notation means}

\paragraph{Complement.}
The notation
\[
S^n \setminus S^m
\]
means the set of all points of \(S^n\) that are not in the chosen embedded copy of \(S^m\).

\paragraph{Homotopy equivalence.}
The notation
\[
X \simeq Y
\]
means that \(X\) and \(Y\) are homotopy equivalent: they may not be homeomorphic, but from the point of view of homotopy theory they have the same essential shape.

Thus
\[
S^n \setminus S^m \simeq S^{\,n-m-1}
\]
means that the complement of \(S^m\) inside \(S^n\) can be continuously deformed onto a sphere of dimension \(n-m-1\).

\subsubsection{Coordinate picture}

Write
\[
\mathbb{R}^{n+1}=\mathbb{R}^{m+1}\times \mathbb{R}^{n-m}.
\]
Then points of \(S^n\subset \mathbb{R}^{n+1}\) can be written as
\[
(u,v),
\qquad
u\in \mathbb{R}^{m+1},
\quad
v\in \mathbb{R}^{n-m},
\quad
\|u\|^2+\|v\|^2=1.
\]

The standard embedded copy of \(S^m\) is
\[
S^m
=
\{(u,0)\in S^n : \|u\|=1\}.
\]
So its complement is
\[
S^n \setminus S^m
=
\{(u,v)\in S^n : v\neq 0\}.
\]

This description is very useful. The condition \(v\neq 0\) means that every point of the complement has a nonzero component in the \(\mathbb{R}^{n-m}\)-direction. Therefore one may keep only the \emph{direction} of \(v\), namely
\[
\frac{v}{\|v\|}\in S^{\,n-m-1}.
\]
This is the geometric source of the sphere
\[
S^{\,n-m-1}.
\]

\subsubsection{Why the dimension \(n-m-1\) appears}

The embedded sphere \(S^m\) lives in the first \(m+1\) coordinates. The remaining coordinates form a space
\[
\mathbb{R}^{n-m}.
\]
A nonzero vector in \(\mathbb{R}^{n-m}\) determines a direction, and the space of directions in \(\mathbb{R}^{n-m}\) is the unit sphere
\[
S^{\,n-m-1}.
\]

So the formula
\[
S^n\setminus S^m \simeq S^{\,n-m-1}
\]
should be understood as saying that after removing \(S^m\), what remains is controlled by the possible \emph{normal directions} away from \(S^m\), and these directions form a sphere of dimension \(n-m-1\).

\subsubsection{The deformation retraction idea}

The proof uses the homotopy
\[
H((u,v),t)
=
\frac{\bigl((1-t)u,v\bigr)}
{\sqrt{(1-t)^2\|u\|^2+\|v\|^2}},
\qquad
0\leq t\leq 1.
\]

This map has the following geometric meaning.

\begin{itemize}
	\item At time \(t=0\), nothing has changed:
	\[
	H((u,v),0)=(u,v).
	\]
	
	\item As \(t\) increases, the \(u\)-part is gradually shrunk toward \(0\), while the \(v\)-part is kept.
	
	\item At time \(t=1\), the point has moved to
	\[
	H((u,v),1)=\left(0,\frac{v}{\|v\|}\right),
	\]
	which lies on the sphere
	\[
	\{(0,w)\in S^n : \|w\|=1\}\cong S^{\,n-m-1}.
	\]
\end{itemize}

Thus every point of the complement slides onto that smaller sphere, and this shows that
\[
S^n\setminus S^m
\]
deformation retracts onto
\[
S^{\,n-m-1}.
\]

\subsubsection{Example 1: \(S^1\setminus S^0\)}

Take
\[
n=1,
\qquad
m=0.
\]
Then the statement becomes
\[
S^1\setminus S^0 \simeq S^{\,1-0-1}=S^0.
\]

Now \(S^0\) is just two points. If we think of \(S^1\) as the circle in \(\mathbb{R}^2\), and \(S^0\subset S^1\) as two opposite points, then removing those two points from the circle leaves two open arcs.

Each open arc is contractible, so the whole space has exactly two connected components, each of which is contractible. Hence the complement has the homotopy type of two points:
\[
S^1\setminus S^0 \simeq S^0.
\]

This is the simplest example of the theorem.

\subsubsection{Example 2: \(S^2\setminus S^0\)}

Now let
\[
n=2,
\qquad
m=0.
\]
Then
\[
S^2\setminus S^0 \simeq S^{\,2-0-1}=S^1.
\]

Geometrically, this means: remove two points from the 2-sphere, for example the north and south poles. What remains is a sphere with two punctures, which is homeomorphic to an open cylinder
\[
S^1\times \mathbb{R}.
\]
An open cylinder deformation retracts onto its circular core \(S^1\). Therefore
\[
S^2\setminus S^0 \simeq S^1.
\]

So the complement of two points in a sphere behaves like a circle.

\subsubsection{Example 3: \(S^2\setminus S^1\)}

Let
\[
n=2,
\qquad
m=1.
\]
Then the formula gives
\[
S^2\setminus S^1 \simeq S^{\,2-1-1}=S^0.
\]

Here \(S^1\subset S^2\) is an equator. Removing the equator divides the sphere into the upper open hemisphere and the lower open hemisphere. Each hemisphere is contractible, so the complement has two contractible connected components. Hence it has the homotopy type of two points:
\[
S^2\setminus S^1 \simeq S^0.
\]

This example shows clearly how the sphere \(S^0\) arises as the set of two ``sides'' of the equator.

\subsubsection{Example 4: \(S^3\setminus S^1\)}

Let
\[
n=3,
\qquad
m=1.
\]
Then
\[
S^3\setminus S^1 \simeq S^{\,3-1-1}=S^1.
\]

So the complement of a standard circle in \(S^3\) has the homotopy type of a circle. This is an important prototype in geometric topology and knot theory: complements of circles and more general knots in \(S^3\) are fundamental objects of study.

Even though \(S^3\setminus S^1\) is a three-dimensional space, its essential homotopy type in this standard case is only that of \(S^1\).

\subsubsection{Example 5: \(S^3\setminus S^2\)}

Let
\[
n=3,
\qquad
m=2.
\]
Then
\[
S^3\setminus S^2 \simeq S^{\,3-2-1}=S^0.
\]

A standard \(S^2\subset S^3\) plays the role of an equator in one higher dimension. Its complement has two connected components, each contractible, so again the complement has the homotopy type of two points.

\subsubsection{A table of low-dimensional cases}

\[
\begin{array}{|c|c|c|c|}
	\hline
	n & m & S^n\setminus S^m & \text{Homotopy type} \\
	\hline
	1 & 0 & S^1\setminus S^0 & S^0 \\
	\hline
	2 & 0 & S^2\setminus S^0 & S^1 \\
	\hline
	2 & 1 & S^2\setminus S^1 & S^0 \\
	\hline
	3 & 0 & S^3\setminus S^0 & S^2 \\
	\hline
	3 & 1 & S^3\setminus S^1 & S^1 \\
	\hline
	3 & 2 & S^3\setminus S^2 & S^0 \\
	\hline
\end{array}
\]

\subsubsection{A concrete coordinate example: the equator in \(S^2\)}

Consider
\[
S^2
=
\{(x,y,z)\in \mathbb{R}^3 : x^2+y^2+z^2=1\}.
\]
Embed
\[
S^1
=
\{(x,y,0)\in \mathbb{R}^3 : x^2+y^2=1\}
\]
as the equator.

Then
\[
S^2\setminus S^1
=
\{(x,y,z)\in S^2 : z\neq 0\}.
\]
This is the union of the upper hemisphere
\[
\{(x,y,z)\in S^2 : z>0\}
\]
and the lower hemisphere
\[
\{(x,y,z)\in S^2 : z<0\}.
\]

Each hemisphere is contractible. Therefore the whole complement has the homotopy type of two points:
\[
S^2\setminus S^1 \simeq S^0.
\]

\subsubsection{What this statement is \emph{not}}

It is important not to confuse this statement with a quotient construction.

A quotient space has the form
\[
X/A,
\]
where the subspace \(A\subset X\) is collapsed to a point. For example,
\[
D^n/\partial D^n \cong S^n.
\]

But in the present theorem we are not collapsing \(S^m\); we are removing it:
\[
S^n\setminus S^m.
\]
So this is a statement about a complement and a homotopy type, not about a quotient space.

\subsubsection{Final geometric interpretation}

The formula
\[
S^n\setminus S^m \simeq S^{\,n-m-1}
\]
should be remembered as follows:

\begin{quote}
	When one removes a standard \(m\)-sphere from an \(n\)-sphere, the essential topology of the complement is determined by the possible normal directions away from that \(m\)-sphere, and those directions form a sphere of dimension \(n-m-1\).
\end{quote}

This is why the complement does not remember the full dimension \(n\), but only the codimension
\[
n-m.
\]
The complement is controlled by the sphere of directions in the normal space, and that sphere is exactly
\[
S^{\,n-m-1}.
\]
	
\end{document}
